% !TeX root = main.tex
\documentclass[11pt]{article}
% Shared preamble of the Guitar Virtual Amp study documents.
% Included by docs/papers/*/main.tex with % Shared preamble of the Guitar Virtual Amp study documents.
% Included by docs/papers/*/main.tex with % Shared preamble of the Guitar Virtual Amp study documents.
% Included by docs/papers/*/main.tex with \input{../common/preamble}.
%
% These are working documents, not submissions: one column, generous
% explanations, and explicit labels for where each claim comes from.

\usepackage[utf8]{inputenc}
\usepackage[T1]{fontenc}
\usepackage{lmodern}
\usepackage[a4paper,margin=2.6cm]{geometry}
\usepackage{microtype}
\usepackage{parskip}
\usepackage{graphicx}
\usepackage{tikz}
\usetikzlibrary{positioning,arrows.meta,fit,backgrounds,calc,shapes.geometric,decorations.pathreplacing}
\usepackage{pgfplots}
\pgfplotsset{compat=1.17}
\usepackage{booktabs}
\usepackage{tabularx}
\usepackage{longtable}
\usepackage{makecell}
\usepackage{multirow}
\usepackage{array}
\usepackage{amsmath,amssymb}
\usepackage[table]{xcolor}
\usepackage{enumitem}
\setlist{itemsep=2pt,topsep=4pt}
\usepackage{caption}
\captionsetup{font=small,labelfont=bf,margin=0.5cm}
\usepackage[numbers,sort&compress]{natbib}
\usepackage[skins,breakable]{tcolorbox}
\usepackage{fancyhdr}
\usepackage[colorlinks=true,linkcolor=blue!50!black,citecolor=green!40!black,urlcolor=blue!60!black]{hyperref}

\setcounter{tocdepth}{2}
\setcounter{secnumdepth}{3}

% ---------------------------------------------------------------- colours
\definecolor{keyc}{HTML}{1F5FAD}
\definecolor{appc}{HTML}{2E7D32}
\definecolor{warnc}{HTML}{C25E00}
\definecolor{examc}{HTML}{555555}
\definecolor{nextc}{HTML}{6A1B9A}

% ---------------------------------------------------------------- boxes
% Intuition before the details.
\newtcolorbox{keyidea}[1][]{enhanced,breakable,colback=keyc!5,colframe=keyc,
  fonttitle=\bfseries,title={Key idea#1},left=6pt,right=6pt,top=4pt,bottom=4pt}
% What Guitar Virtual Amp does, with code pointers.
\newtcolorbox{inapp}[1][]{enhanced,breakable,colback=appc!5,colframe=appc,
  fonttitle=\bfseries,title={In Guitar Virtual Amp#1},left=6pt,right=6pt,top=4pt,bottom=4pt}
% Pitfalls, caveats, things that are easy to misread.
\newtcolorbox{caution}[1][]{enhanced,breakable,colback=warnc!5,colframe=warnc,
  fonttitle=\bfseries,title={Caution#1},left=6pt,right=6pt,top=4pt,bottom=4pt}
% A small calculation done by hand.
\newtcolorbox{worked}[1][]{enhanced,breakable,colback=examc!4,colframe=examc,
  fonttitle=\bfseries,title={Worked example#1},left=6pt,right=6pt,top=4pt,bottom=4pt}
% Planned or in-progress work.
\newtcolorbox{nextsteps}[1][]{enhanced,breakable,colback=nextc!4,colframe=nextc,
  fonttitle=\bfseries,title={Next steps#1},left=6pt,right=6pt,top=4pt,bottom=4pt}

% ---------------------------------------------------------------- evidence labels
% Every number or claim about a method carries one of these.
\DeclareRobustCommand{\evtag}[2]{{\setlength{\fboxsep}{1.3pt}%
  \colorbox{#1!15}{\textcolor{#1!70!black}{\scriptsize\sffamily\bfseries #2}}}}
\DeclareRobustCommand{\measured}{\evtag{appc}{MEASURED}}  % run by us, in this repository
\DeclareRobustCommand{\reported}{\evtag{keyc}{REPORTED}}  % stated by the authors
\DeclareRobustCommand{\expected}{\evtag{warnc}{EXPECTED}} % hypothesis, not yet measured
% Status of a component of the app.
\DeclareRobustCommand{\done}{\evtag{appc}{BUILT}}
\DeclareRobustCommand{\wip}{\evtag{warnc}{IN PROGRESS}}
\DeclareRobustCommand{\planned}{\evtag{nextc}{PLANNED}}

% ---------------------------------------------------------------- shorthands
% Underscores and slashes need no escaping inside these.
\newcommand{\code}[1]{{\small\nolinkurl{#1}}}
\newcommand{\file}[1]{{\small\nolinkurl{#1}}}
% Shell commands and flags with spaces (\code drops spaces).
\newcommand{\cmd}[1]{\texttt{\small\detokenize{#1}}}
\newcommand{\ms}{\,ms}
\newcommand{\dB}{\,dB}
\newcommand{\Hz}{\,Hz}
\newcommand{\kHz}{\,kHz}
\newcommand{\term}[1]{\emph{#1}}

\pagestyle{fancy}
\fancyhf{}
\renewcommand{\headrulewidth}{0.3pt}
\fancyfoot[C]{\thepage}
.
%
% These are working documents, not submissions: one column, generous
% explanations, and explicit labels for where each claim comes from.

\usepackage[utf8]{inputenc}
\usepackage[T1]{fontenc}
\usepackage{lmodern}
\usepackage[a4paper,margin=2.6cm]{geometry}
\usepackage{microtype}
\usepackage{parskip}
\usepackage{graphicx}
\usepackage{tikz}
\usetikzlibrary{positioning,arrows.meta,fit,backgrounds,calc,shapes.geometric,decorations.pathreplacing}
\usepackage{pgfplots}
\pgfplotsset{compat=1.17}
\usepackage{booktabs}
\usepackage{tabularx}
\usepackage{longtable}
\usepackage{makecell}
\usepackage{multirow}
\usepackage{array}
\usepackage{amsmath,amssymb}
\usepackage[table]{xcolor}
\usepackage{enumitem}
\setlist{itemsep=2pt,topsep=4pt}
\usepackage{caption}
\captionsetup{font=small,labelfont=bf,margin=0.5cm}
\usepackage[numbers,sort&compress]{natbib}
\usepackage[skins,breakable]{tcolorbox}
\usepackage{fancyhdr}
\usepackage[colorlinks=true,linkcolor=blue!50!black,citecolor=green!40!black,urlcolor=blue!60!black]{hyperref}

\setcounter{tocdepth}{2}
\setcounter{secnumdepth}{3}

% ---------------------------------------------------------------- colours
\definecolor{keyc}{HTML}{1F5FAD}
\definecolor{appc}{HTML}{2E7D32}
\definecolor{warnc}{HTML}{C25E00}
\definecolor{examc}{HTML}{555555}
\definecolor{nextc}{HTML}{6A1B9A}

% ---------------------------------------------------------------- boxes
% Intuition before the details.
\newtcolorbox{keyidea}[1][]{enhanced,breakable,colback=keyc!5,colframe=keyc,
  fonttitle=\bfseries,title={Key idea#1},left=6pt,right=6pt,top=4pt,bottom=4pt}
% What Guitar Virtual Amp does, with code pointers.
\newtcolorbox{inapp}[1][]{enhanced,breakable,colback=appc!5,colframe=appc,
  fonttitle=\bfseries,title={In Guitar Virtual Amp#1},left=6pt,right=6pt,top=4pt,bottom=4pt}
% Pitfalls, caveats, things that are easy to misread.
\newtcolorbox{caution}[1][]{enhanced,breakable,colback=warnc!5,colframe=warnc,
  fonttitle=\bfseries,title={Caution#1},left=6pt,right=6pt,top=4pt,bottom=4pt}
% A small calculation done by hand.
\newtcolorbox{worked}[1][]{enhanced,breakable,colback=examc!4,colframe=examc,
  fonttitle=\bfseries,title={Worked example#1},left=6pt,right=6pt,top=4pt,bottom=4pt}
% Planned or in-progress work.
\newtcolorbox{nextsteps}[1][]{enhanced,breakable,colback=nextc!4,colframe=nextc,
  fonttitle=\bfseries,title={Next steps#1},left=6pt,right=6pt,top=4pt,bottom=4pt}

% ---------------------------------------------------------------- evidence labels
% Every number or claim about a method carries one of these.
\DeclareRobustCommand{\evtag}[2]{{\setlength{\fboxsep}{1.3pt}%
  \colorbox{#1!15}{\textcolor{#1!70!black}{\scriptsize\sffamily\bfseries #2}}}}
\DeclareRobustCommand{\measured}{\evtag{appc}{MEASURED}}  % run by us, in this repository
\DeclareRobustCommand{\reported}{\evtag{keyc}{REPORTED}}  % stated by the authors
\DeclareRobustCommand{\expected}{\evtag{warnc}{EXPECTED}} % hypothesis, not yet measured
% Status of a component of the app.
\DeclareRobustCommand{\done}{\evtag{appc}{BUILT}}
\DeclareRobustCommand{\wip}{\evtag{warnc}{IN PROGRESS}}
\DeclareRobustCommand{\planned}{\evtag{nextc}{PLANNED}}

% ---------------------------------------------------------------- shorthands
% Underscores and slashes need no escaping inside these.
\newcommand{\code}[1]{{\small\nolinkurl{#1}}}
\newcommand{\file}[1]{{\small\nolinkurl{#1}}}
% Shell commands and flags with spaces (\code drops spaces).
\newcommand{\cmd}[1]{\texttt{\small\detokenize{#1}}}
\newcommand{\ms}{\,ms}
\newcommand{\dB}{\,dB}
\newcommand{\Hz}{\,Hz}
\newcommand{\kHz}{\,kHz}
\newcommand{\term}[1]{\emph{#1}}

\pagestyle{fancy}
\fancyhf{}
\renewcommand{\headrulewidth}{0.3pt}
\fancyfoot[C]{\thepage}
.
%
% These are working documents, not submissions: one column, generous
% explanations, and explicit labels for where each claim comes from.

\usepackage[utf8]{inputenc}
\usepackage[T1]{fontenc}
\usepackage{lmodern}
\usepackage[a4paper,margin=2.6cm]{geometry}
\usepackage{microtype}
\usepackage{parskip}
\usepackage{graphicx}
\usepackage{tikz}
\usetikzlibrary{positioning,arrows.meta,fit,backgrounds,calc,shapes.geometric,decorations.pathreplacing}
\usepackage{pgfplots}
\pgfplotsset{compat=1.17}
\usepackage{booktabs}
\usepackage{tabularx}
\usepackage{longtable}
\usepackage{makecell}
\usepackage{multirow}
\usepackage{array}
\usepackage{amsmath,amssymb}
\usepackage[table]{xcolor}
\usepackage{enumitem}
\setlist{itemsep=2pt,topsep=4pt}
\usepackage{caption}
\captionsetup{font=small,labelfont=bf,margin=0.5cm}
\usepackage[numbers,sort&compress]{natbib}
\usepackage[skins,breakable]{tcolorbox}
\usepackage{fancyhdr}
\usepackage[colorlinks=true,linkcolor=blue!50!black,citecolor=green!40!black,urlcolor=blue!60!black]{hyperref}

\setcounter{tocdepth}{2}
\setcounter{secnumdepth}{3}

% ---------------------------------------------------------------- colours
\definecolor{keyc}{HTML}{1F5FAD}
\definecolor{appc}{HTML}{2E7D32}
\definecolor{warnc}{HTML}{C25E00}
\definecolor{examc}{HTML}{555555}
\definecolor{nextc}{HTML}{6A1B9A}

% ---------------------------------------------------------------- boxes
% Intuition before the details.
\newtcolorbox{keyidea}[1][]{enhanced,breakable,colback=keyc!5,colframe=keyc,
  fonttitle=\bfseries,title={Key idea#1},left=6pt,right=6pt,top=4pt,bottom=4pt}
% What Guitar Virtual Amp does, with code pointers.
\newtcolorbox{inapp}[1][]{enhanced,breakable,colback=appc!5,colframe=appc,
  fonttitle=\bfseries,title={In Guitar Virtual Amp#1},left=6pt,right=6pt,top=4pt,bottom=4pt}
% Pitfalls, caveats, things that are easy to misread.
\newtcolorbox{caution}[1][]{enhanced,breakable,colback=warnc!5,colframe=warnc,
  fonttitle=\bfseries,title={Caution#1},left=6pt,right=6pt,top=4pt,bottom=4pt}
% A small calculation done by hand.
\newtcolorbox{worked}[1][]{enhanced,breakable,colback=examc!4,colframe=examc,
  fonttitle=\bfseries,title={Worked example#1},left=6pt,right=6pt,top=4pt,bottom=4pt}
% Planned or in-progress work.
\newtcolorbox{nextsteps}[1][]{enhanced,breakable,colback=nextc!4,colframe=nextc,
  fonttitle=\bfseries,title={Next steps#1},left=6pt,right=6pt,top=4pt,bottom=4pt}

% ---------------------------------------------------------------- evidence labels
% Every number or claim about a method carries one of these.
\DeclareRobustCommand{\evtag}[2]{{\setlength{\fboxsep}{1.3pt}%
  \colorbox{#1!15}{\textcolor{#1!70!black}{\scriptsize\sffamily\bfseries #2}}}}
\DeclareRobustCommand{\measured}{\evtag{appc}{MEASURED}}  % run by us, in this repository
\DeclareRobustCommand{\reported}{\evtag{keyc}{REPORTED}}  % stated by the authors
\DeclareRobustCommand{\expected}{\evtag{warnc}{EXPECTED}} % hypothesis, not yet measured
% Status of a component of the app.
\DeclareRobustCommand{\done}{\evtag{appc}{BUILT}}
\DeclareRobustCommand{\wip}{\evtag{warnc}{IN PROGRESS}}
\DeclareRobustCommand{\planned}{\evtag{nextc}{PLANNED}}

% ---------------------------------------------------------------- shorthands
% Underscores and slashes need no escaping inside these.
\newcommand{\code}[1]{{\small\nolinkurl{#1}}}
\newcommand{\file}[1]{{\small\nolinkurl{#1}}}
% Shell commands and flags with spaces (\code drops spaces).
\newcommand{\cmd}[1]{\texttt{\small\detokenize{#1}}}
\newcommand{\ms}{\,ms}
\newcommand{\dB}{\,dB}
\newcommand{\Hz}{\,Hz}
\newcommand{\kHz}{\,kHz}
\newcommand{\term}[1]{\emph{#1}}

\pagestyle{fancy}
\fancyhf{}
\renewcommand{\headrulewidth}{0.3pt}
\fancyfoot[C]{\thepage}


\fancyhead[L]{\small Tablature transcription notes}
\fancyhead[R]{\small\nouppercase{\leftmark}}

\title{Automatic Guitar Tablature Transcription\\[4pt]
\large State of the art, how it works, and what Guitar Virtual Amp does}
\author{Breno de Angelo\\\small Guitar Virtual Amp project notes}
\date{Snapshot of 25 November 2026}

\begin{document}
\maketitle

\begin{abstract}
\noindent
This document explains, from the ground up, how a computer can turn a song
recording into a guitar tablature, what the published research has achieved at
each step, and exactly what Guitar Virtual Amp (the app in this repository)
does at each step today. It is written to be read by the project owner, not by
reviewers: every technical term is defined before it is used, every number is
labelled as \measured{} (run by us, in this repository), \reported{} (stated by
the authors of a paper) or \expected{} (a hypothesis not yet tested), and every
part of the app is labelled as \done{}, \wip{} or \planned{}. The short version:
the practical state of the art for a commercial song is a \emph{modular
pipeline} (separate the guitar, transcribe its notes, find beats and bars,
choose strings and frets, mark techniques). The app implements one. On our
benchmarks its weakest links are distorted electric guitar inside real mixes
(note F1 about 0.56 on distorted guitars of real songs) and, before the learned
fingering model, the choice of string on melodic lines.
\end{abstract}

\tableofcontents
\clearpage

% !TeX root = ../main.tex
\section*{How to read this document}
\addcontentsline{toc}{section}{How to read this document}

\paragraph{Purpose.} Coding agents work on the transcription pipeline of Guitar
Virtual Amp, run benchmarks and change models. This document is the place where
their work is explained in one coherent story: what problem each piece solves,
what the literature knows about it, what the app does, how well it works, and
what is planned. It replaces an earlier review written in conference-paper
style. It is deliberately longer and more explicit than a paper.

\paragraph{Structure.}
\begin{itemize}
\item \textbf{Part I, Background} (Sections~\ref{sec:problem}--\ref{sec:eval})
  defines the task and teaches the minimum of audio processing, machine
  learning and evaluation needed to follow the rest. Skip it if you know these.
\item \textbf{Part II, The pipeline} (Sections~\ref{sec:datasets}--\ref{sec:e2e})
  walks through the stages in the order audio flows through them. Each section
  follows the same template: \emph{what the stage must do}, \emph{how methods
  work}, \emph{state of the art}, an ``In Guitar Virtual Amp'' box, and
  \emph{failure modes}.
\item \textbf{Part III, Our system} (Sections~\ref{sec:app}--\ref{sec:assessment})
  shows the whole app pipeline, all our benchmark results with their protocol,
  and an assessment: where we stand against the state of the art, which
  decisions were taken and why, and what the agents are working on next.
\item \textbf{Appendices}: a glossary and a map from source files to sections,
  with the commands that reproduce each benchmark.
\end{itemize}

\paragraph{Labels.} The coloured boxes have fixed meanings:
\begin{keyidea}[]
The intuition behind a method, before its details.
\end{keyidea}
\begin{inapp}[]
What the app does, with file names under \file{guitar_core/src/} unless another
path is given, and the status of each part.
\end{inapp}
\begin{caution}[]
A pitfall: a number that is easy to misread, a comparison that is not fair, a
known gap.
\end{caution}
Inline tags say where a claim comes from: \measured{} means we ran it in this
repository (the script is named), \reported{} means a paper states it (with a
citation), and \expected{} means it is a hypothesis that no one has measured
yet. For parts of the app: \done{} exists in the code, \wip{} is being worked
on, \planned{} is only in the roadmap.

\paragraph{Related files.} \file{ROADMAP.md} (project status),
\file{docs/transcription_roadmap.md} (the detailed research log, from which most
of the measured numbers here are taken), \file{CLAUDE.md} (architecture notes for
agents), \file{scripts/train/README.md} (training and Python benchmarks). The
companion document on estimating a recording's guitar tone is
\file{docs/papers/tone-estimation/}.

\paragraph{Snapshot.} Numbers and code descriptions are as of 25 November 2026
(commit \code{25afb6d} plus the uncommitted work described where relevant).
When the code changes, the ``In Guitar Virtual Amp'' boxes and Part~III are the
parts to update.

\subsection*{The whole story in one table}
\addcontentsline{toc}{subsection}{The whole story in one table}

Table~\ref{tab:summary} summarises the rest of the document: for each stage,
the best published approach, what the app uses, our best measured number and
the main remaining gap. Every entry is explained later.

\begin{table}[h]
\centering
\caption{Stage-by-stage summary. F50 = note F1 with a 50\ms{} onset window
(Section~\ref{sec:eval}). ``v2'' = our electric-guitar fine-tune of the GAPS
model. All numbers in the last two columns are \measured{} unless marked.}
\label{tab:summary}
\small
\begin{tabularx}{\textwidth}{@{}>{\raggedright}p{2.0cm}>{\raggedright}p{3.1cm}>{\raggedright}p{3.0cm}>{\raggedright}p{3.4cm}X@{}}
\toprule
Stage & Best published & App uses & Our best number & Main gap \\
\midrule
Separation (\S\ref{sec:separation}) & Band-split transformers (BS-RoFormer) & BS-RoFormer SW, else Demucs \code{htdemucs_6s} & Real songs: F50 0.41 (mix) $\to$ 0.70 (stem) & All guitars in one stem \\
Note transcription (\S\ref{sec:notes}) & Guitar CRNNs (GAPS); MuScriptor on mixes & GAPS CRNN; v2 as local option & EGSet12 (real amp) 0.89; GOAT amp renders 0.71 & Distorted guitar in real mixes: 0.56 \\
Rhythm (\S\ref{sec:rhythm}) & Beat This! & Beat This! final0 + bar Viterbi & Beat F 0.97 (accompaniment), 0.55 (solo) & Solo lines; compound meters \\
Fingering (\S\ref{sec:fingering}) & Learned seq2seq, audio-conditioned (TART) & Learned tagger (local), else Viterbi & String accuracy, true notes: 0.62 solo, 0.81 EGDB & Few players in training data \\
Techniques (\S\ref{sec:techniques}) & Learned per-note classifier (TART) & Rules on the pitch contour & Bends found 35--56\%, vibrato 20--40\% & No palm mute, harmonics; distortion \\
Whole tab (\S\ref{sec:benchmark}) & TART: 54\% Tab F1, solo guitar \reported{} & All of the above & Tab F1 0.65 on separated mixes (tagger) & No real-song benchmark yet \\
\bottomrule
\end{tabularx}
\end{table}


\part{Background}
% !TeX root = ../main.tex
\section{The problem}
\label{sec:problem}

\subsection{What a tablature is}

A \term{tablature} (tab) is guitar notation that says \emph{where} to play each
note, not only \emph{which} note. It has one horizontal line per string, drawn
with the highest-pitched string (string 1, the thin high E) at the top. A number
on a line is the fret to press on that string; 0 means the open string. Numbers
stacked vertically are played together (a chord). Above or below the lines, a
tab also gives the rhythm (how long each note lasts, where the bars are) and
technique marks: \emph{b} for a bend, \emph{/} or \emph{\textbackslash} for a
slide, \emph{h} and \emph{p} for hammer-on and pull-off, \emph{\textasciitilde}
for vibrato, \emph{PM} for palm mute. Figure~\ref{fig:tab} shows a short example.

\begin{figure}[h]
\centering
\begin{tikzpicture}[x=0.9cm,y=0.42cm,font=\small]
  % strings
  \foreach \s/\n in {1/e,2/B,3/G,4/D,5/A,6/E} {
    \draw[gray] (0,-\s) -- (13,-\s);
    \node[anchor=east,font=\scriptsize] at (-0.1,-\s) {\n};
  }
  \draw (0,-1) -- (0,-6); \draw (6.5,-1) -- (6.5,-6); \draw (13,-1) -- (13,-6);
  % bar 1: a power chord twice, then a single note line
  \foreach \x in {0.8,2.0} {
    \node[fill=white,inner sep=1pt] at (\x,-6) {3};
    \node[fill=white,inner sep=1pt] at (\x,-5) {5};
    \node[fill=white,inner sep=1pt] at (\x,-4) {5};
  }
  \node[fill=white,inner sep=1pt] at (3.3,-3) {5};
  \node[fill=white,inner sep=1pt] at (4.2,-3) {7};
  \node[font=\scriptsize] at (3.75,-2.3) {h};
  \node[fill=white,inner sep=1pt] at (5.4,-2) {8};
  \node[font=\scriptsize] at (5.4,-0.2) {b (full)};
  \draw[->,>=Stealth] (5.55,-1.6) to[bend left=30] (5.9,-0.6);
  % bar 2
  \node[fill=white,inner sep=1pt] at (7.5,-1) {5};
  \node[fill=white,inner sep=1pt] at (8.8,-1) {7};
  \node[font=\scriptsize] at (8.15,-0.3) {/};
  \node[fill=white,inner sep=1pt] at (10.2,-1) {7};
  \node[font=\scriptsize] at (10.6,-0.3) {\textasciitilde\textasciitilde\textasciitilde};
  \node[fill=white,inner sep=1pt] at (12,-6) {0};
  \node[font=\scriptsize] at (12,-7) {PM};
  % rhythm hints
  \foreach \x in {0.8,2.0,3.3,4.2,5.4,7.5,8.8,10.2,12} \draw (\x,-7.8) -- (\x,-8.8);
  \node[font=\scriptsize,anchor=west] at (0,-9.6) {rhythm stems and beams go here; bar lines split the time into measures};
\end{tikzpicture}
\caption{A two-bar tablature. Bar 1: a power chord (frets 3--5--5 on strings
6--5--4) struck twice, a hammer-on from fret 5 to 7 on string 3, then fret 8 on
string 2 bent up a whole tone. Bar 2: a slide from fret 5 to 7 on string 1, a
held note with vibrato, and a palm-muted open low E.}
\label{fig:tab}
\end{figure}

A complete tablature of one guitar part therefore contains four kinds of
information:
\begin{enumerate}
\item \textbf{Notes}: when each note starts (\term{onset}), when it stops
  (\term{offset}) and its pitch.
\item \textbf{Fingering}: a string and fret for each note, given a tuning
  (standard E A D G B E, drop D, a half step down, 7 strings, \ldots) and maybe
  a capo.
\item \textbf{Rhythm notation}: beats, bars, time signatures (4/4, 6/8, \ldots)
  and a written duration for each note (quarter, eighth, triplet, \ldots).
\item \textbf{Techniques}: bends (and how far), releases, slides, hammer-ons
  and pull-offs, vibrato, palm mutes, harmonics, dead notes.
\end{enumerate}
When a song has several guitars, each part must also become its own track.

\subsection{Why it is hard}

\paragraph{One pitch, several places.} On a guitar the same pitch can be played
on several strings. The E above middle C (MIDI pitch 64) is the open first
string, the 5th fret of string 2, the 9th fret of string 3, the 14th of string 4
and the 19th of string 5: five positions on a 22-fret guitar
(Figure~\ref{fig:positions}). The audio says the pitch; the tab needs the
position. Guitarists choose positions so the hand moves little and chords are
playable, and they have habits, so ``the'' correct tab is partly a matter of
convention. This is the \term{fingering} or \term{string-fret assignment}
problem (Section~\ref{sec:fingering}).

\begin{figure}[h]
\centering
\begin{tikzpicture}[x=0.55cm,y=0.42cm,font=\scriptsize]
  \foreach \s in {1,...,6} \draw[gray] (0,-\s) -- (22,-\s);
  \foreach \f in {0,...,22} \draw[gray!50] (\f,-1) -- (\f,-6);
  \draw[very thick] (0,-1) -- (0,-6);
  \foreach \f in {3,5,7,9,15,17,19,21} \fill[gray!40] (\f-0.5,-3.5) circle (2pt);
  \fill[gray!40] (11.5,-2.5) circle (2pt); \fill[gray!40] (11.5,-4.5) circle (2pt);
  \foreach \s/\f in {1/0,2/5,3/9,4/14,5/19} {
    \fill[keyc] ({max(\f-0.5,-0.4)},-\s) circle (5pt);
    \node[text=white,font=\tiny\bfseries] at ({max(\f-0.5,-0.4)},-\s) {\f};
  }
  \foreach \s/\n in {1/1 (e),2/2 (B),3/3 (G),4/4 (D),5/5 (A),6/6 (E)}
    \node[anchor=east] at (-0.6,-\s) {string \n};
  \foreach \f in {0,5,10,15,20} \node at (\f,-6.9) {\f};
\end{tikzpicture}
\caption{The five places to play E4 (MIDI 64) on a 22-fret guitar in standard
tuning. Numbers are frets.}
\label{fig:positions}
\end{figure}

\paragraph{Polyphony.} A strummed chord puts up to six notes within a few tens
of milliseconds, and their harmonics overlap. Detecting which pitches are
present is much harder than for a single melody.

\paragraph{Distortion.} Electric guitar is almost always played through an
amplifier and effects. Distortion adds many harmonics and also
\term{intermodulation} products (sums and differences of the frequencies being
played). A note detector can then report notes that are not there, typically an
octave or a fifth away from the real ones (Section~\ref{sec:audio-harmonics}).

\paragraph{The mix.} In a commercial song the guitar shares the recording with
drums, bass, keyboards and voice, often with a second guitar playing the same
or a different part. The guitar must be isolated before (or while) its notes
are found (Section~\ref{sec:separation}).

\paragraph{Scarce labels.} To train or test a system we need recordings where
every note's onset, pitch and string is known. That requires special hardware
(a \term{hexaphonic pickup}, which records each string separately) or hours of
expert annotation. The public datasets are small and mostly solo guitar
(Section~\ref{sec:datasets}). No public dataset pairs commercial band
recordings with reference tabs.

\subsection{Two ways to build a system}

\begin{keyidea}[: modular versus end-to-end]
A \textbf{modular pipeline} splits the task into sub-problems (separation,
notes, rhythm, fingering, techniques) and solves each with its own model or
algorithm. Each part can be tested and replaced on its own, and it can start
from a full mix. Errors accumulate from stage to stage. An \textbf{end-to-end
model} predicts strings and frets directly from audio with one network. It can
use the timbre of each string, but all published ones need an isolated guitar.
\end{keyidea}

Figure~\ref{fig:families} contrasts the two. A third line of work, large
multi-instrument \term{sequence-to-sequence} transcribers (MT3, YourMT3+,
MuScriptor, Section~\ref{sec:notes-seq2seq}), does not output tablature but is a
strong candidate for the note stage of a modular pipeline, because it is
trained on full mixes.

\begin{figure}[h]
\centering
\begin{tikzpicture}[
  node distance=4mm,
  box/.style={draw, rounded corners=2pt, font=\small, align=center,
              minimum height=8mm, inner sep=3pt, fill=white},
  arr/.style={-{Stealth[length=2mm]}, semithick}]
  \node[box] (mix) {Song\\(mix)};
  \node[box, right=of mix] (sep) {Source\\separation};
  \node[box, right=of sep] (amt) {Note\\transcription};
  \node[box, right=of amt] (fing) {String and\\fret};
  \node[box, right=of fing] (tab) {Tab};
  \node[box, above=of amt] (beat) {Beats and bars};
  \node[box, below=of amt] (tech) {Techniques};
  \draw[arr] (mix) -- (sep);
  \draw[arr] (sep) -- (amt);
  \draw[arr] (amt) -- (fing);
  \draw[arr] (fing) -- (tab);
  \draw[arr] (mix.north) |- (beat.west);
  \draw[arr] (beat.east) -| (tab.north);
  \draw[arr] (sep.south) |- (tech.west);
  \draw[arr] (tech.east) -| (tab.south);
  \begin{scope}[on background layer]
    \node[draw, dashed, rounded corners, fit=(mix)(tab)(beat)(tech),
          inner sep=3mm, label={[font=\small]above left:Modular pipeline}] {};
  \end{scope}
  \node[box, below=20mm of mix] (mix2) {Solo\\guitar};
  \node[box, right=8mm of mix2, minimum width=60mm] (e2e)
       {One network: audio $\to$ fret on each string};
  \node[box, right=8mm of e2e] (tab2) {Tab};
  \draw[arr] (mix2) -- (e2e);
  \draw[arr] (e2e) -- (tab2);
  \begin{scope}[on background layer]
    \node[draw, dashed, rounded corners, fit=(mix2)(tab2), inner sep=3mm,
          label={[font=\small]below left:End-to-end model}] {};
  \end{scope}
\end{tikzpicture}
\caption{The two families of tablature transcription systems. Modular
pipelines can start from a full mix; every published end-to-end model starts
from an isolated guitar recording.}
\label{fig:families}
\end{figure}

\begin{inapp}[: which family]
The app is a modular pipeline (\file{transcribe/mod.rs}, function
\code{transcribe::run}). Section~\ref{sec:app} shows it end to end; Part~II
explains each stage. The reason for this choice is measured, not assumed: on
electric guitar and on mixes the modular pipeline beats the end-to-end models
we could run (Section~\ref{sec:bench-tab}).
\end{inapp}

% !TeX root = ../main.tex
\section{Audio primer}
\label{sec:audio}

This section introduces the signal-processing ideas used everywhere else. It
assumes nothing beyond high-school mathematics.

\subsection{Samples, frames and the spectrogram}

A digital recording is a list of numbers, the \term{samples}, measured at a
fixed \term{sample rate}. The app's audio engine runs at 48\,000 samples per
second (48\kHz); the models resample to what they were trained on (16\kHz for
the guitar CRNN, 22.05\kHz for Basic Pitch and Beat This!, 44.1\kHz stereo for
BS-RoFormer).

Pitch is easier to see in frequency than in time. The \term{short-time Fourier
transform} (STFT) cuts the signal into overlapping \term{frames} of $N$ samples
(the \term{window}, e.g.\ 2048), moves forward by a \term{hop} of $H$ samples
(e.g.\ 512) between frames, and computes the spectrum of each frame: how much
energy there is at each frequency. Stacking these spectra side by side gives a
\term{spectrogram}, an image with time on the horizontal axis, frequency on the
vertical axis and energy as brightness. Most transcription models take a
spectrogram as input.

\begin{worked}[: resolution]
At 16\kHz, a hop of 160 samples gives $16000/160 = 100$ frames per second, so
time is resolved in steps of 10\ms. A window of $N$ samples resolves frequency
in steps of $f_s/N$: with $N=2048$ at 16\kHz that is 7.8\Hz, less than a
semitone even at the lowest guitar note (E2 = 82.4\Hz; the next semitone, F2, is
87.3\Hz). Longer windows give finer frequency but blur time. This trade-off is
why some systems compute two resolutions.
\end{worked}

Three frequency axes are common:
\begin{itemize}
\item \textbf{Linear} (the plain STFT): equal spacing in Hz.
\item \textbf{Mel}: spacing that follows human pitch perception, roughly linear
  below 1\kHz and logarithmic above. Beat This! uses 128 mel bands.
\item \textbf{Constant-Q transform} (CQT): bins spaced by a fixed musical
  interval (for example a third of a semitone), so every octave has the same
  number of bins. Natural for music; TabCNN and Basic Pitch use CQT-like inputs.
\end{itemize}
Energies are almost always compressed with a logarithm (decibels), because
loudness perception is roughly logarithmic and it keeps quiet notes visible.

\subsection{Pitch, MIDI numbers and harmonics}
\label{sec:audio-harmonics}

Musical pitch is labelled here with \term{MIDI numbers}: A4 = 440\Hz{} is 69,
and each semitone adds one. The frequency of MIDI pitch $p$ is
\[
  f(p) = 440 \cdot 2^{(p-69)/12}\ \text{Hz}.
\]
Standard guitar tuning is 40 (E2), 45, 50, 55, 59, 64 (E4) from string 6 to
string 1; each fret adds one semitone.

A plucked string does not vibrate at one frequency. It produces a
\term{fundamental} $f_0$ plus \term{harmonics} at $2f_0, 3f_0, 4f_0,\ldots$ The
2nd harmonic is one octave up, the 3rd an octave and a fifth, the 4th two
octaves. Every note detector must decide whether energy at $2f_0$ is a harmonic
of a note at $f_0$ or a separate note an octave higher. Getting this wrong
produces an \term{octave error}.

\begin{keyidea}[: why distortion causes octave-low errors]
A distortion circuit clips the waveform. Clipping a sum of two notes at
frequencies $f_1$ and $f_2$ creates energy at $f_2 - f_1$, $2f_1 - f_2$ and
similar combinations (\term{intermodulation}). For a power chord (a root and the
fifth above it, ratio 3:2) the difference tone $f_2 - f_1 = f_1/2$ is one
octave \emph{below} the root. A model trained on clean guitar sees a strong
partial an octave below and reports a note there. This is the main failure we
measure on the ``Metal High Gain'' preset (Section~\ref{sec:bench-app}): most
octave errors are an octave low.
\end{keyidea}

\subsection{Onsets and the pitch contour}

An \term{onset} is the moment a note starts: a burst of broadband energy from
the pick, followed by the harmonic tone. Onset detectors look for sudden rises
in energy across frequencies (\term{spectral flux}: the sum of positive
frame-to-frame increases of the spectrum). A \term{pitch contour} is the
fundamental frequency over time at fine resolution (a fraction of a semitone).
A bend appears in it as a smooth rise of one or two semitones, vibrato as a
periodic wobble of a quarter to half a semitone at 4--8\Hz, a slide as a
staircase of semitone steps. Section~\ref{sec:techniques} builds technique
detection on this.

\subsection{Stereo and panning}

A stereo recording has a left and a right channel. Mixing engineers
\term{pan} each instrument: a guitar panned hard left is only in the left
channel, one in the centre is equal in both. Rock and metal records often place
two rhythm guitars hard left and hard right and a lead guitar in the centre. The
ratio of left to right energy in a time-frequency bin therefore tells roughly
where in the stereo image the sound came from. The app uses this to split
guitars that a separator returns as one stem (Section~\ref{sec:pan-split}).

% !TeX root = ../main.tex
\section{Machine learning and decoding primer}
\label{sec:learning}

\subsection{Supervised learning in one page}

A \term{model} is a function with many adjustable numbers (\term{parameters} or
\term{weights}). \term{Training} adjusts them so that, on example inputs, the
model's outputs approach known correct answers (\term{labels}), by minimising a
\term{loss} that measures the disagreement. For transcription, the input is a
spectrogram and the labels say, for each frame and pitch, whether a note starts
or sounds there.

Data is split three ways. The \term{training set} is what the model learns
from. The \term{validation set} is used to make choices during development
(when to stop training, which settings to keep). The \term{test set} is used
only to report the final number; if any choice is made by looking at it, the
number becomes optimistic. A test set is \term{held out} when no part of it, nor
the same performers, appears in training. A result is \term{zero-shot} when the
model is tested on a dataset it never saw at all, which is the situation of a
user's song.

\begin{caution}[: overfitting and leakage]
A model that scores far better on its training players than on new ones has
\term{overfit}: it memorised instead of generalising. Our fingering tagger
scores 0.95 string accuracy on its training players and 0.57--0.69 on new
ones (Section~\ref{sec:fingering-app}). \term{Leakage} is subtler: if test
recordings share performers, songs or rendered tones with the training data,
the test score is inflated. The rule in this project is that GuitarSet players
03--05 are never used to fit anything, so they stay a clean test of fingering.
\end{caution}

\term{Fine-tuning} starts from a model already trained on one dataset and
continues training on another, so a small new dataset is enough. Our electric
guitar model is the GAPS classical-guitar model fine-tuned on electric guitar.
\term{Data augmentation} creates extra training examples by transforming the
ones we have; here, playing clean DI recordings through many amplifier
settings so the model sees the same notes in many tones.

\subsection{Network families mentioned in this document}

\begin{description}[style=nextline,leftmargin=1.2cm]
\item[CNN (convolutional neural network)] Slides small learned filters over the
  spectrogram image; good at local patterns such as a note's harmonic stack.
\item[RNN / GRU / LSTM (recurrent networks)] Read a sequence one step at a time
  and carry a memory; good at how notes evolve over frames.
\item[CRNN] A CNN followed by a recurrent layer. The GAPS guitar model (Kong et
  al.'s architecture) is a CRNN.
\item[Transformer] Relates every position of a sequence to every other one
  through \term{attention}~\cite{vaswani2017attention}. Used by BS-RoFormer,
  Beat This!, MT3-style transcribers and our fingering tagger.
\item[Encoder-decoder / sequence-to-sequence] An encoder reads the input; a
  decoder writes an output sequence one \term{token} at a time, each token
  conditioned on the ones before (\term{autoregressive} decoding). MT3,
  YourMT3+ and MuScriptor write notes as tokens such as ``time 1.23 s'',
  ``instrument: distorted guitar'', ``note on 52''. Decoding one token at a time
  is slow on a CPU.
\item[Diffusion model] Learns to turn random noise into a sample step by step;
  Noise2Fret uses one to generate tablature.
\end{description}

\subsection{From network output to notes}

Most note models output a \term{posteriorgram}: for each frame and each of the
88 piano pitches, a number between 0 and 1 that the model believes is the
probability of an onset (or of a sounding note) there. Turning it into a list of
notes requires \term{post-processing}: find peaks above a threshold in the
onset map, follow the note while the frame map stays above another threshold,
drop notes that are too short. These thresholds are not learned and change
results noticeably; in the app they live next to each model
(\file{transcribe/guitar_crnn.rs}, \file{transcribe/basic_pitch.rs}).

\begin{inapp}[: how models run]
All neural models used inside the app are exported to ONNX (a standard file
format for networks) and run in pure Rust with the \code{tract-onnx} crate
(version 0.23), on the CPU. Basic Pitch is compiled into the binary; the guitar
CRNN and Beat This! are optional downloads; the fingering tagger and the
electric CRNN are local files the user copies in. Source separators are the
exception: they run as external Python programs (MSST or Demucs) in their own
virtual environments, on the GPU when PyTorch finds one.
\end{inapp}

\subsection{Dynamic programming and the Viterbi algorithm}
\label{sec:viterbi}

Several stages must pick the best \emph{sequence} of choices: a string for each
note, a subdivision for each beat, a bar length for each bar. Choosing each
item on its own ignores how consecutive choices fit together (moving the hand
from fret 2 to fret 14 between two fast notes is bad even if each position is
fine alone). Trying every sequence is impossible: 20 notes with 3 positions each
give $3^{20} \approx 3.5$ billion sequences.

\begin{keyidea}[: Viterbi decoding]
Define a cost for each choice on its own (\term{local cost}) and a cost for each
pair of consecutive choices (\term{transition cost}). The best sequence ending in
choice $j$ at step $t$ is the best sequence ending in some choice $i$ at step
$t-1$, plus the transition $i\to j$, plus the local cost of $j$. So we only need
to keep, for each choice at each step, the cost of the best path that reaches it
and a pointer to where it came from. At the end we follow the pointers back.
The work is (steps)~$\times$~(choices)$^2$ instead of
(choices)$^{\text{steps}}$~\cite{viterbi1967error}.
\end{keyidea}

\begin{worked}[: two notes, three positions]
Note 1 can be played at fret 0, 5 or 10; note 2 at fret 2, 7 or 12. Local cost:
0.1 per fret (prefer low positions). Transition cost: 0.3 per fret of hand
movement. Best costs after note 1: $\{0, 0.5, 1.0\}$. For note 2 at fret 7:
from fret 0 costs $0 + 0.3\cdot 7 + 0.7 = 2.8$; from fret 5 costs $0.5 + 0.3\cdot
2 + 0.7 = 1.8$; from fret 10 costs $1.0 + 0.9 + 0.7 = 2.6$. Keep 1.8 with a
pointer to fret 5. Doing the same for frets 2 and 12 and taking the cheapest
end gives fret 0 then fret 2 (total $0 + 0.6 + 0.2 = 0.8$). If note 1 had been
forced high, the transition cost would have pulled note 2 high too: that is
the hand-position memory the method provides.
\end{worked}

A \term{hidden Markov model} (HMM) is the probabilistic version of the same
idea: costs are negative log-probabilities. A \term{dynamic Bayesian network}
(DBN) generalises it to several hidden variables (for beats: tempo and position
in the bar). Our fingering, our beat subdivision choice and our bar decoding are
all Viterbi searches with hand-set costs.

% !TeX root = ../main.tex
\section{How results are measured}
\label{sec:eval}

Numbers are only comparable when they use the same metric, tolerance and test
set. This section defines every metric used later and lists which of them each
of our benchmarks reports.

\subsection{Precision, recall and F1 for notes}

Compare a list of predicted notes with a list of reference (correct) notes. A
prediction \term{matches} a reference note when the pitches are equal (within
half a semitone) and the onsets are within a tolerance $\tau$. Each reference
note can be matched at most once (a one-to-one matching). Then
\[
  \text{precision} = \frac{\text{matched}}{\text{predicted}},\qquad
  \text{recall} = \frac{\text{matched}}{\text{reference}},\qquad
  F_1 = \frac{2\cdot\text{precision}\cdot\text{recall}}{\text{precision}+\text{recall}}.
\]
Precision drops when the system invents notes; recall drops when it misses
notes. F1 is high only when both are high.

\begin{worked}[: F1]
A bar has 10 reference notes. The system outputs 12 notes, 8 of which match.
Precision $= 8/12 = 0.67$, recall $= 8/10 = 0.80$,
$F_1 = 2\cdot0.67\cdot0.80/(0.67+0.80) = 0.73$.
\end{worked}

Variants used in this document:
\begin{description}[style=nextline,leftmargin=1.2cm]
\item[F50 (onset F1)] $\tau = 50\ms$, the \code{mir_eval}
  convention~\cite{raffel2014mireval,hawthorne2018onsets}. Computed on the raw
  notes of a model, before any rhythm quantisation. This is what papers report
  and what our Python stage benchmarks use.
\item[Onset-offset F1] Also requires the predicted offset to be within
  $\max(50\ms, 0.2\,d)$ of the reference offset ($d$: the reference duration).
  Measures note lengths.
\item[Final-tab F1 (ours)] $\tau = 80\ms$, computed by the Rust benchmarks on the
  notes of the finished tab, \emph{after} the notes were snapped to the rhythm
  grid. Quantisation moves onsets by up to half a grid step, hence the wider
  tolerance. Reported as the \emph{median} over items.
\item[Tab F1] A note match that is also on the same string (and hence fret) as
  the reference. Measures the whole pipeline.
\item[String accuracy] Given the \emph{correct} notes, the share placed on the
  reference string. Measures the fingering stage alone.
\item[Octave errors] The share of reference notes that were found only one or two
  octaves away. Diagnoses distortion problems.
\end{description}

\subsection{Rhythm metrics}

Beat trackers are scored with a beat \term{F-measure} (a predicted beat is
correct within 70\ms{} of a reference beat) and continuity metrics such as
\term{CMLt}, which counts only runs of consecutive correct beats, so a tracker
that keeps losing the beat scores low even if many beats are individually
right~\cite{foscarin2024beatthis}. Downbeat (first beat of each bar)
F-measure is computed the same way. Our final-tab benchmarks also count songs
whose tempo is right and whose bar lines fall on the reference bar lines.

\subsection{Separation metrics}

Separators are usually scored by the \term{signal-to-distortion ratio} (SDR),
in decibels: the energy of the reference stem divided by the energy of the
error, so higher is better and 3\dB{} more means half the error
energy~\cite{vincent2006bss}. SDR rewards a stem that \emph{sounds} close to the
reference; it does not say whether the stem is easy to transcribe. We therefore
score separators by the F50 of a transcription of their stem.

\subsection{Which number is which}

Table~\ref{tab:which-number} lists where each kind of number in this document
comes from.

\begin{table}[h]
\centering
\caption{Metrics produced by each benchmark in the repository.}
\label{tab:which-number}
\small
\begin{tabularx}{\textwidth}{@{}>{\raggedright}p{4.6cm}>{\raggedright}p{3.6cm}X@{}}
\toprule
Benchmark (code) & Metric & Notes \\
\midrule
\code{scripts/train/bench_notes.py}, \code{bench_separation.py}, \code{bench_tab.py}, \code{bench_beats.py} & F50, onset-offset F1, Tab F1, beat F / CMLt & Python, \code{mir_eval}; raw model notes; mean per item; up to 12 items per group \\
\code{bench_moisesdb.py} & F50 against notes of the isolated guitars & No human labels: ranks separators, not note models (Section~\ref{sec:bench-moises}) \\
\code{eval_guitarset} and other ignored tests in \file{transcribe/eval.rs} & Final-tab F1 (80\ms), string/fret agreement, tempo and bar lines & Rust, the real app pipeline end to end; median per item \\
\code{eval_fingering}, \code{examples/fingering.rs} & String accuracy on true notes & Fingering alone \\
\bottomrule
\end{tabularx}
\end{table}

\begin{caution}[: comparing numbers]
\begin{itemize}
\item An F50 of 0.80 and a final-tab F1 of 0.80 are not the same thing: one is
  on raw notes with a 50\ms{} window, the other on the finished tab with an
  80\ms{} window.
\item Mean and median differ when a few items fail badly; Python benchmarks
  report means, Rust benchmarks medians.
\item Labels carry timing offsets. Guitar-TECHS labels (from a MIDI pickup) are
  25--40\ms{} late, IDMT labels about 30\ms{} early, GOAT up to $\pm40\ms$
  \measured{} (\code{align_labels.py}). With a 50\ms{} tolerance this alone can
  move F1 a lot, so our benchmarks shift each recording's labels by its median
  offset first.
\item Many published results are on the dataset the model was trained on
  (different files, same players or studio). The use case is zero-shot.
\item Agreement with one human tab is a lower bound on quality: several
  fingerings of the same passage can be equally good.
\end{itemize}
\end{caution}


\part{The pipeline, stage by stage}
% !TeX root = ../main.tex
\section{Datasets}
\label{sec:datasets}

\subsection{What makes a dataset useful here}

For transcribing a guitar part from a commercial song, four properties decide
how useful a dataset is: (i) whether the audio is \emph{electric} guitar with
realistic amplification, (ii) whether labels include the \emph{string} (not just
the pitch) and the \emph{techniques}, (iii) whether the guitar appears inside a
\emph{mix}, and (iv) the \emph{licence}, which decides whether a model trained
on it can be distributed. No public dataset has all four.

A recurring term is \term{DI} (direct input): the clean signal straight from the
guitar's pickup, before any amplifier. A DI recording can be played through any
number of amplifier simulations afterwards (\term{re-amping} or
\term{rendering}), which turns one labelled performance into many labelled
tones. This is how EGDB, GOAT and our own training data multiply their size.

\subsection{The datasets}

\paragraph{Audio with string labels.}
GuitarSet~\cite{xi2018guitarset} records acoustic guitar with a hexaphonic
pickup: per-string notes and pitch contours for 360 excerpts by six players
(00--05), each excerpt played as accompaniment (``comp'') and as a melody
(``solo''). It is the standard benchmark, and its small size makes overfitting a
recurring concern~\cite{zang2024synthtab}.
EGDB~\cite{chen2022egdb} provides 240 tablatures played on electric guitar with
onset, pitch and string labels, as DI and through five amplifier simulations,
plus a few real recordings; EGDB-PG~\cite{chen2025egdbpg} re-renders it through
many amplifier-cabinet configurations.
GOAT~\cite{loth2025goat} has 5.9\,h of DI electric guitar with Guitar Pro tabs
that include techniques, and 29.5\,h of amplifier renders.
EGSet12~\cite{pedroza2024tones} has 12 real performances through a miked
amplifier with string labels.
IDMT-SMT-Guitar~\cite{kehling2014idmt} has isolated notes and licks labelled
with expression styles (bend, slide, vibrato, harmonics, dead notes).

\paragraph{Audio with note labels only.}
GAPS~\cite{riley2024gaps} aligns published classical guitar scores to 14\,h of
recordings by more than 200 performers: the largest real guitar set with note
labels. Guitar-TECHS~\cite{pedroza2025guitartechs} records electric guitar at
once as DI, miked amplifier and room microphones, with MIDI labels from a
hexaphonic MIDI pickup, covering techniques, chords, scales and musical
excerpts.

\paragraph{Symbolic and synthetic.}
DadaGP~\cite{sarmento2021dadagp} is 26\,181 Guitar Pro songs (no audio), the
main resource for learning fingering from human tabs; it is available for
research on request. SynthTab~\cite{zang2024synthtab} renders DadaGP through
commercial guitar plugins. GuitarFlow~\cite{loth2025guitarflow} renders tabs
with a sampler and then transforms the result toward realistic electric guitar.

\paragraph{Mixes.}
MUSDB18~\cite{rafii2017musdb18} (150 songs) and
MoisesDB~\cite{pereira2023moisesdb} (240 songs, with guitar sub-stems split into
acoustic, clean electric and distorted electric) are separation datasets: real
multitrack songs, no note labels. MuScriptor~\cite{rouard2026muscriptor} built a
private set of 170\,000 real recordings (over 11\,000\,h) with notes aligned by
audio-to-score synchronisation~\cite{muller2021sync}.

\begin{table}[h]
\centering
\caption{Datasets and how this project uses them. Licences as recorded in
\file{scripts/train/README.md} and \file{docs/transcription_roadmap.md}; NC =
non-commercial, ND = no derivatives.}
\label{tab:datasets}
\small
\begin{tabularx}{\textwidth}{@{}>{\raggedright}p{2.3cm}>{\raggedright}p{3.0cm}ccc>{\raggedright}p{2.0cm}X@{}}
\toprule
Dataset & Audio & String & Tech. & Mix & Licence & Our use \\
\midrule
GuitarSet & acoustic, mic + hexaphonic & yes & contours & no & CC BY 4.0 & Players 00--02 train; 03--05 test (never fitted on) \\
EGDB & electric DI + 5 amp sims & yes & no & no & none stated & Train: ids not multiple of 5; test: multiples of 5 and the real clips \\
EGSet12 & real electric, miked amp & yes & no & no & CC BY 4.0 & Test only \\
GOAT & electric DI + amp renders & yes (tabs) & yes & no & NC, on request & Official test split; never in a distributed model \\
Guitar-TECHS & electric DI, amp, mics & no & yes & no & CC BY 4.0 & Chords/scales train; excerpts and techniques test \\
IDMT-SMT-Guitar & electric notes and licks & yes & yes & no & CC BY-NC-ND & Test only (techniques) \\
GAPS & classical, real & no & no & no & -- & Pre-training of the GAPS model (by its authors) \\
DadaGP & none (symbolic) & yes & yes & -- & research, on request & \planned{} fingering pre-training \\
MUSDB18 & real stems & no & no & yes & NC research & Backing for synthetic band mixes \\
MoisesDB & real stems & no & no & yes & check terms & Real-mix benchmark (Section~\ref{sec:bench-moises}) \\
\bottomrule
\end{tabularx}
\end{table}

\begin{inapp}[: where the data lives]
None of the datasets is in the repository. They live under
\file{~/.cache/gva_test} (about 60\,GB), prepared by the scripts in
\file{scripts/train/} (\code{prepare_data.py}, \code{make_mixes.py},
\code{prepare_moisesdb.py}, \code{align_labels.py}), which write manifests
(\code{.tsv} files listing audio and notes files) that the Rust benchmark
\code{eval_manifest} and the Python benchmarks read.
\end{inapp}

\begin{caution}[: the missing dataset]
There is no public set of commercial band recordings with reference tabs, which
is exactly the use case. MuScriptor showed how one could be built (align
existing tabs to recordings, reject bad alignments). The app has the tool for
this, \code{scripts/train/align_tabs.py} \done, but no aligned and hand-checked
song pairs yet (Section~\ref{sec:bench-real}).
\end{caution}

% !TeX root = ../main.tex
\section{Source separation}
\label{sec:separation}

\subsection{What the stage must do}

Input: the stereo song. Output: a \term{stem} per instrument, here at least
guitar, bass and drums. The guitar stem feeds the note model; the drum stem can
help the beat tracker; the bass stem gets its own transcription.

\subsection{How separators work}

\begin{keyidea}[: masking]
Most separators compute the mix's spectrogram and predict, for each
time-frequency bin, which fraction of the energy belongs to the target
instrument (a \term{mask} between 0 and 1). Multiplying the mix spectrogram by
the mask and inverting the STFT gives the stem. The network learns masks from
songs whose individual stems are known (MUSDB18, MoisesDB and private sets).
\end{keyidea}

Two families matter here.
\begin{itemize}
\item \textbf{Hybrid Transformer Demucs} (HT Demucs)~\cite{rouard2023htdemucs}
  processes the waveform and the spectrogram in two U-Net branches joined by a
  transformer. It reports 9.20\dB{} SDR on MUSDB18 with extra training data
  \reported. Its six-source variant \code{htdemucs_6s} adds guitar and piano
  stems and is the most widely available open guitar separator (MIT licence).
\item \textbf{Band-split models} cut the complex spectrogram into frequency
  bands and model each band and the relations between bands with transformers.
  BS-RoFormer~\cite{lu2023bsroformer} won the music track of the Sound Demixing
  Challenge 2023 with 9.80\dB{} SDR on MUSDB18HQ without extra data
  \reported; Mel-RoFormer~\cite{wang2023melroformer} uses overlapping mel bands
  instead. The MSST framework~\cite{solovyev2026msst} trains and runs these
  models, and community checkpoints with a guitar stem are distributed through
  it.
\end{itemize}
A third idea, \textbf{query-based separation}
(Banquet~\cite{watcharasupat2024banquet}), uses one decoder that is told which
instrument to extract. With 24.9\,M parameters it beat \code{htdemucs_6s} on
guitar and piano on MoisesDB \reported{} and can target narrow classes such as
clean acoustic guitar.

\subsection{Two limitations for tablature}

\paragraph{One stem for all guitars.} Separators return one guitar stem, so a
rhythm guitar and a lead guitar, or two double-tracked rhythm guitars, come out
together. Their notes then mix in one transcription.

\paragraph{Artefacts.} Separated stems contain \term{bleed} (traces of
keyboards, synth pads or vocals) and smeared attacks. Note models are trained on
clean, isolated guitar, so separation artefacts are inputs they have never
seen. They become false or missed onsets.

\subsection{Splitting guitars by stereo position}
\label{sec:pan-split}

Spatial cues can undo part of the first limitation. \term{ADRess} (azimuth
discrimination and resynthesis)~\cite{barry2004adress} assigns each
time-frequency bin of a stereo signal to a position between hard left and hard
right from the ratio of its left and right magnitudes. Keeping only the bins
near one position extracts what was panned there.

\begin{inapp}[: separation and pan split \done]
\begin{itemize}
\item \textbf{Separator} (\file{transcribe/separation.rs}): BS-RoFormer SW (a
  community six-stem checkpoint, run through MSST by the script
  \code{separate_msst.py} in its own Python environment, installed by
  \code{scripts/install_bs_roformer.sh}) when installed; otherwise Demucs
  \code{htdemucs_6s} (\code{scripts/install_demucs.sh}).
  \code{GVA_SEPARATOR=demucs|bs-roformer} forces one. Both retry on the CPU if
  the GPU runs out of memory. The Rust side decodes and resamples the song to
  44.1\kHz{} stereo for BS-RoFormer. On a 261\,s song on the laptop GPU (RTX
  4070, 8\,GB): BS-RoFormer 34\,s and about 3.5\,GB of GPU memory, Demucs 10\,s
  \measured. Without a separator the app transcribes the full mix.
\item \textbf{Pan split} (\file{transcribe/stereo.rs}): the stereo guitar stem
  is split into left, centre and right parts with soft masks on a 2048/512 STFT
  (centre: $|\text{pan}| < 0.2$; sides: $|\text{pan}| > 0.5$). A part with at
  least 12\% of the stem's energy is transcribed on its own. If its notes agree
  with a louder part (note F1 $\geq 0.5$, allowing an echo lag of up to
  0.8\,s), it is the same guitar (one guitar panned halfway, a double-tracked
  part or a stereo delay) and is dropped; so are parts with fewer than a tenth
  of the most notes. The remaining parts become tracks, the most chord-heavy
  first (``Guitarra base'', ``Guitarra solo''). \code{GVA_SPLIT_GUITARS=0} turns
  it off.
\end{itemize}
\end{inapp}

\subsection{What we measured}

Separators are scored here by what matters downstream: the F50 of a note
transcription of their guitar stem (Section~\ref{sec:eval}).

\begin{table}[h]
\centering
\caption{F50 of guitar notes in 41 synthetic band mixes (labelled guitar over
MUSDB18 drums, bass and vocals), transcribed from the raw mix or from a
separated stem \measured{} (\code{bench_separation.py}). The same guitar parts
without backing score 0.82 with the v2 model.}
\label{tab:sep-synth}
\small
\begin{tabular}{@{}lccc@{}}
\toprule
Note model & Raw mix & \code{htdemucs_6s} stem & BS-RoFormer stem \\
\midrule
GAPS CRNN & 0.51 & 0.72 & \textbf{0.81} \\
v2 CRNN (electric fine-tune) & 0.50 & 0.74 & \textbf{0.81} \\
Basic Pitch & 0.36 & 0.58 & 0.64 \\
YourMT3+ & 0.51 & 0.45 & 0.52 \\
MuScriptor large & 0.44 & 0.64 & 0.75 \\
\bottomrule
\end{tabular}
\end{table}

On real songs (MoisesDB, 222 songs, no human note labels; protocol in
Section~\ref{sec:bench-moises}), the v2 CRNN goes from 0.41 F50 on the raw mix
to 0.58 on the \code{htdemucs_6s} stem and 0.70 on the BS-RoFormer stem
\measured. Separation is the largest single gain in the pipeline, and the gap
between the two separators is larger on real songs (0.12) than on our synthetic
mixes (0.07).

For the pan split \measured{} (\code{eval_guitar_pairs}, GuitarSet accompaniment
and solo mixed into one stereo file, 11 mixes per layout): with rhythm hard left
and lead hard right the app outputs two tracks in 10 of 11 mixes and lead-part
F1 rises from 0.40 (one merged track) to 0.85; with rhythm left and lead centre,
11 of 11 and 0.43 to 0.88. A single guitar panned halfway or with a 300\ms{}
stereo echo stays one track in 11 of 11, with unchanged F1.

\subsection{Failure modes}
\begin{itemize}
\item Two guitars in the same position of the stereo image stay one track
  (known, by construction).
\item Keyboards and synth pads can leak into the guitar stem \expected.
\item BS-RoFormer SW is about 3.4 times slower than Demucs on the laptop GPU and
  needs about 4\,GB of free GPU memory \measured. When the GPU is full it falls
  back to the CPU, which is slower.
\item The BS-RoFormer SW checkpoint has no stated licence (it comes from a
  mirror of a deleted account): it is downloaded for local use only and must
  not be bundled.
\end{itemize}

% !TeX root = ../main.tex
\section{Note transcription}
\label{sec:notes}

\subsection{What the stage must do}

Input: the guitar stem (or the mix). Output: a list of notes, each with onset,
offset, pitch and a loudness estimate. This is \term{automatic music
transcription} (AMT) restricted to one instrument.

\subsection{Frame-and-onset models}

\begin{keyidea}[: predict onsets and sustain separately]
Onsets and Frames~\cite{hawthorne2018onsets} trains one network to output two
posteriorgrams: where notes \emph{start} and where notes \emph{sound}. A note is
created only at a detected onset and lasts while the frame output stays high.
Separating the two makes repeated notes and sustained notes both readable.
\end{keyidea}

Kong et al.~\cite{kong2021highres} improved this for piano by \emph{regressing}
the exact onset and offset time inside each 10\ms{} frame instead of a yes/no
target, which gives sub-frame timing and tolerates labels that are a little off.
It reached 96.72\% onset F1 on the MAESTRO piano dataset \reported. This
architecture (a CRNN with onset, offset, frame and velocity outputs at 100
frames per second) is the basis of today's best guitar models: Riley et al.\
first adapted it to guitar by aligning published scores to
recordings~\cite{riley2024domain}, then trained it on the GAPS classical guitar
dataset~\cite{riley2024gaps}, with state-of-the-art results on GuitarSet both
when trained on it and zero-shot \reported.

\textbf{Basic Pitch}~\cite{bittner2022basicpitch} is a small model that works
on any instrument. It outputs onsets, notes and a \term{pitch contour} at a
third of a semitone. It is fast and a useful baseline, and its contour is the
basis of our technique detection, but being instrument-agnostic it tends to
report string harmonics as extra notes.

\subsection{Electric guitar and tone}

Amplifier tone is the main obstacle. EGDB showed that accuracy depends strongly
on the amplifier~\cite{chen2022egdb}. EGDB-PG trained a Tone-informed
Transformer (TIT) on many amplifier-cabinet renders with a tone
embedding~\cite{chen2025egdbpg}; Pedroza et al.\ trained on synthetic
performances played through recorded tones and effects~\cite{pedroza2024tones}.
The recipe is the same in all of them: show the model the same performance in
many tones.

\subsection{Sequence-to-sequence transcribers}
\label{sec:notes-seq2seq}

Hawthorne et al.~\cite{hawthorne2021seq2seq} showed that a generic transformer
that ``translates'' a spectrogram into MIDI-like tokens matches specialised piano
models. MT3~\cite{gardner2022mt3} extended this to many instruments with tokens
that include the instrument. YourMT3+~\cite{chang2024yourmt3} improved the
encoder and the training mixtures but remains limited on pop recordings.
MuScriptor~\cite{rouard2026muscriptor} changes the data rather than the
architecture: a decoder-only transformer (up to 1.3\,B parameters, 5\,s
inputs) pre-trained on 1.45\,M synthesised MIDI files, trained on 170\,000
aligned real recordings, then refined with reinforcement learning. Distorted
electric guitar is the second most frequent class in its real data (about 47\%
of tracks). It can be told which instruments are present. Two limitations
matter for us: its tokens cannot represent the same pitch twice at once on one
instrument (keeping such notes in its test set lowers the large model's onset
F1 from 60.4 to 51.8 \reported), and like all models of this kind it outputs no
string and merges all guitars of one class. Its weights are non-commercial.

\begin{inapp}[: note models \done]
\begin{itemize}
\item \textbf{Guitar, default}: the GAPS CRNN (\file{transcribe/guitar_crnn.rs},
  optional download by \code{scripts/install_transcription_models.sh}, exported
  to ONNX by \code{scripts/export_guitar_crnn.py}, MIT). Audio is resampled to
  16\kHz{} and cut into 10\,s windows with 50\% overlap, run in parallel on at
  most 4 threads (each window takes about 600\,MB in tract). Notes are read
  with Kong's post-processing: onset and offset peaks with sub-frame shift, a
  frame threshold, notes under 30\ms{} dropped.
\item \textbf{Guitar, electric option}: \code{guitar_electric_crnn.onnx}, our
  fine-tune (``v2'', below), chosen in the Transcribe form or with
  \cmd{--guitar-model electric}. Local file only, because EGDB has no licence.
\item \textbf{Fallback and bass}: Basic Pitch (\file{transcribe/basic_pitch.rs},
  bundled, Apache-2.0; 22.05\kHz, 2\,s windows). It always transcribes the bass,
  and the guitar when the CRNN is not installed. With the CRNN it still runs on
  the guitar audio, only for its pitch contour (Section~\ref{sec:techniques}).
\item \code{prune_weak} removes weak simultaneous detections and, for Basic
  Pitch, likely overtones.
\end{itemize}
\end{inapp}

\begin{inapp}[: how v2 was trained]
\code{scripts/train/finetune_crnn.py} starts from the GAPS weights and trains on
EGDB (training ids), Guitar-TECHS chords and scales and GuitarSet players
00--02. Each DI recording is rendered through the app's own factory presets with
the desktop tool \code{render_tones} (\file{backend/offline.rs}); ``Metal High
Gain'' and ``Blues Crunch'' are left out of training because they are the
distortion benchmark. Labels are first shifted per recording to remove timing
offsets (\code{align_labels.py}). BatchNorm layers are frozen (batches are
small), precision is fp32/TF32 (fp16 overflowed in the power spectrum; bf16
crashed with NaNs at step 1100). The adopted run (\code{runs/spark_tf32}) took
8000 steps of 8 clips, about 1\,h on one DGX Spark; the best validation
checkpoint was step 5000.
\end{inapp}

\subsection{What we measured}

\begin{table}[h]
\centering
\caption{Note F50 (mean per item) on held-out sets \measured{}
(\code{bench_notes.py}, up to 12 items per group). MS = MuScriptor, conditioned
on the guitar classes; multi-instrument models are scored on their guitar
notes. Speed in audio seconds per second on one DGX Spark GPU (Basic Pitch on
CPU).}
\label{tab:bench-notes}
\footnotesize
\setlength{\tabcolsep}{4pt}
\begin{tabular}{@{}lccccccc@{}}
\toprule
Set & GAPS & v2 & Basic Pitch & YourMT3+ & MS small & MS medium & MS large \\
\midrule
EGDB, DI + 3 amp sims & 0.74 & \textbf{0.84} & 0.67 & 0.71 & 0.73 & 0.74 & 0.76 \\
EGDB real recordings (5) & 0.89 & \textbf{0.90} & 0.69 & 0.83 & 0.71 & 0.79 & 0.59 \\
EGSet12, real, miked amp & \textbf{0.89} & 0.88 & 0.76 & 0.54 & 0.85 & 0.88 & \textbf{0.89} \\
GOAT DI & 0.82 & \textbf{0.84} & 0.74 & 0.81 & 0.74 & 0.83 & 0.82 \\
GOAT, 5 amp renders & 0.56 & 0.71 & 0.56 & 0.39 & 0.66 & 0.76 & \textbf{0.77} \\
Guitar-TECHS excerpts & \textbf{0.81} & 0.78 & 0.61 & 0.70 & 0.66 & 0.73 & 0.71 \\
GuitarSet comp & 0.87 & 0.89 & 0.75 & \textbf{0.92} & 0.77 & 0.77 & 0.77 \\
GuitarSet solo & 0.91 & 0.92 & 0.81 & \textbf{0.93} & 0.86 & 0.89 & 0.86 \\
\midrule
Speed ($\times$ real time) & 34 & 89 & 165 & 14 & 6.4 & 4.0 & 1.3 \\
\bottomrule
\end{tabular}
\end{table}

How to read Table~\ref{tab:bench-notes}:
\begin{itemize}
\item \textbf{No model wins everywhere.} On amplified guitar the models trained
  on many tones lead: MuScriptor large 0.77 and v2 0.71 on the GOAT amplifier
  renders, against 0.56 for GAPS. MuScriptor large also makes the fewest
  octave-low errors there (2.1\% of notes against 8.0\% for GAPS).
\item \textbf{Real electric guitar through a clean amp} (EGSet12): GAPS, v2 and
  MuScriptor large tie at 0.88--0.89.
\item \textbf{YourMT3+} is best on acoustic GuitarSet and has the best note
  ends, but labels much amplified guitar as other instruments (0.39 on GOAT
  amps).
\item \textbf{Speed}: MuScriptor large runs at 1.3$\times$ real time on a large
  GPU, too slow for the app's CPU path; v2 runs at 89$\times$.
\end{itemize}

In the full app pipeline (\code{eval_guitarset}, GuitarSet players 03--05 played
through the app's presets, final-tab F1 median / octave errors) \measured:
\begin{center}
\small
\begin{tabular}{@{}lcc@{}}
\toprule
Tone & GAPS & v2 (electric) \\
\midrule
Clean & 0.84 / 1.2\% & \textbf{0.86} / 0.8\% \\
``Blues Crunch'' & 0.76 / 1.5\% & \textbf{0.81} / 1.0\% \\
``Metal High Gain'' & 0.43 / 4.5\% & \textbf{0.69} / 2.3\% \\
\bottomrule
\end{tabular}
\end{center}
A hand-written rule that raised notes struck with a strong partial a twelfth
above them (the distortion difference tone) made clean and crunch worse and
metal no better, so it was dropped: distortion needs a model trained on
distorted guitar, which is what v2 is.

\subsection{Failure modes}
\begin{itemize}
\item \textbf{Distortion}: octave-low errors from difference tones
  (Section~\ref{sec:audio-harmonics}). The 8-string SoundFont render of
  ``Meshuggah -- Bleed'' gets a final-tab F1 of 0.02 with the GAPS CRNN, against
  0.24 with Basic Pitch \measured.
\item \textbf{Distorted guitar in real mixes}: on MoisesDB, v2 on the
  BS-RoFormer stem scores 0.56 on distorted guitars, 0.75 on clean electric and
  0.83 on acoustic \measured. This is the largest remaining loss on real songs.
\item \textbf{Separation artefacts} are unseen in training (no published guitar
  transcriber is trained on separated stems); planned as training task T1
  (Section~\ref{sec:next}).
\item \textbf{Same-pitch notes} (unisons on two strings) cannot be represented
  by MT3-style tokens.
\end{itemize}

% !TeX root = ../main.tex
\section{Beats, bars and rhythm notation}
\label{sec:rhythm}

\subsection{What the stage must do}

A tab needs bars, a time signature per bar and a written duration per note.
That takes three steps: find the \term{beats} (the pulse a listener taps along
to), find the \term{downbeats} (the first beat of each bar, which fixes the
meter: 4 beats per bar for 4/4, 3 for 3/4), and \term{quantise} each note onto
a grid of beat subdivisions so it can be written as a quarter, eighth,
sixteenth or triplet.

\subsection{How beat trackers work}

\begin{keyidea}[: activation, then decoding]
A tracker first computes, for every short frame, how likely a beat is there (an
\term{activation}): classically from an onset function such as spectral flux,
today from a neural network. It then \emph{decodes} a regular sequence of beats
from the activation with dynamic programming (Section~\ref{sec:viterbi}), using
the prior knowledge that beats are roughly evenly spaced.
\end{keyidea}

Ellis's tracker~\cite{ellis2007beat} picks the tempo by autocorrelation of the
onset function and then the beat times by dynamic programming. Neural trackers
predict beat and downbeat activations and decode them with a dynamic Bayesian
network that assumes a fixed meter and a tempo range~\cite{bock2016joint};
All-In-One~\cite{kim2023allinone} also tracks song sections.
\textbf{Beat This!}~\cite{foscarin2024beatthis} drops the DBN: a transformer
trained on diverse data (including solo instruments and time-signature changes)
predicts beats and downbeats that are only peak-picked, and surpasses prior
work in F1 while being weaker on continuity \reported.
AlignBeat~\cite{ahn2025alignbeat} reports 3.6 points more downbeat CMLt than
Beat This! \reported; Explore This!~\cite{kihlborg2026explore} predicts a finer
\term{tatum} grid and groups it into beats and bars. A failure analysis on the
SMC dataset~\cite{ahn2026smc} found that current trackers are ``confident but
wrong'' on expressive music, and that the usual DBN minimum tempo of 55\,BPM
forces double-tempo output on 21\% of its tracks \reported.

\subsection{From beats to written rhythm}

Quantisation and the choice of notated durations are rarely studied in the
tablature literature; TART uses a rule-based tab generator~\cite{gupta2026tart}.
The difficulty is that every beat can be divided in two ways that matter in
popular music: \term{straight} (two eighths, four sixteenths) or
\term{triplet} (three eighth-note triplets). A note 1/3 of a beat late is a
triplet in a shuffle and a timing error in a straight rock song.

\begin{inapp}[: rhythm \done]
\begin{itemize}
\item \textbf{Beat This!} (\file{transcribe/beat_net.rs}; ONNX files
  \code{beat_this_final0.onnx}, or \code{beat_this_small0.onnx} when final0 is
  missing, plus \code{beat_this_logmel.onnx}): runs on the mix at 22.05\kHz{},
  128 log-mel bands, 50 activations per second, in 30\,s chunks. Peaks become
  beats; \code{regularize} fills missed beats and drops doubled ones.
\item \textbf{Bars and meter} (\code{beat_net::decode_downbeats}): a
  \term{bar-pointer} Viterbi whose states are ``beat $k$ of a bar of $n$
  beats'', $n$ from 2 to 7. Each meter has a cost (4/4 is favoured,
  \code{METERS}) and changing meter costs \code{METER_CHANGE}; downbeat
  activations are scored against their mean over the song. So a song can switch
  meter or insert a single 2/4 bar.
\item \textbf{Fallback} (\file{transcribe/beats.rs}) when Beat This! is not
  installed: spectral flux, autocorrelation tempo with a 120\,BPM prior, Ellis
  dynamic programming, downbeat phase from the low-frequency band; 4/4 only. It
  runs on the drum stem when there is one.
\item \textbf{Grid} (\code{beats::bar_grid}, \code{arrange::BeatGrid}): whole bars
  around the notes, one length per bar.
\item \textbf{Quantisation} (\file{transcribe/arrange.rs}): onsets within 35\ms{}
  of each other form a chord; \code{choose_subdivisions} picks, for each beat,
  sixteenths or eighth-note triplets by a Viterbi search with a prior towards
  straight time; positions are stored in ticks (960 per beat). \code{notate}
  writes tuplets only inside partly used triplet beats and ties long notes
  across beats; each bar gets its own tempo in the score.
\end{itemize}
\end{inapp}

\subsection{What we measured}

\begin{table}[h]
\centering
\caption{Beat tracking on GuitarSet players 03--05, all 180 excerpts
\measured{} (\code{bench_beats.py}, \code{mir_eval}).}
\label{tab:bench-beats}
\small
\begin{tabular}{@{}lcc@{}}
\toprule
 & Accompaniment (comp) & Solo \\
Tracker & beat F / CMLt / downbeat F & beat F / CMLt / downbeat F \\
\midrule
Beat This! small0 & 0.92 / 0.80 / 0.92 & 0.51 / 0.36 / 0.33 \\
Beat This! final0 (app default) & \textbf{0.97 / 0.92 / 0.94} & \textbf{0.55} / 0.35 / \textbf{0.36} \\
Ellis DP (librosa) & 0.54 / 0.42 / -- & 0.39 / 0.26 / -- \\
\bottomrule
\end{tabular}
\end{table}

\begin{itemize}
\item final0 is a clear gain on accompaniment (CMLt 0.80 to 0.92) for 8\,MB more
  weights.
\item Every tracker fails on unaccompanied solo lines (beat F about 0.5): a
  single melody carries little beat information. In a song the tracker runs on
  the full mix, so this matters mostly for solo-guitar recordings.
\item In the app pipeline (\code{eval_guitarset}, players 03--05, 60 excerpts,
  GAPS), final0 against small0: tempo right in 51 against 49 excerpts, bar lines
  aligned in 41 against 35 (solo 11 against 6, comp 30 against 29) \measured.
  The beat stage takes 13.5\,s against 8.3\,s on a 261\,s song.
\end{itemize}

\begin{caution}[: two sets of rhythm numbers]
The README and an older ROADMAP table report ``tempo right in 59 of 60, bar
lines aligned in 56--57 of 60'' on GuitarSet. Those came from an earlier run
with a different selection of 60 excerpts (and small0); the numbers above are
on players 03--05 only. Even small0 scores only 49 of 60 on players 03--05, so
the selection, not the model, explains most of the gap, but nobody has checked
which excerpts differ. Quote the players 03--05 numbers as current.
\end{caution}

\subsection{Failure modes}
\begin{itemize}
\item Compound meters (6/8, 12/8) are written as 2/4 or 6/4 with triplets
  \measured.
\item Sixteenth-note triplets and 32nd notes cannot be written (the grid has
  only sixteenths and eighth triplets).
\item Slow, rubato or expressive music: confident but wrong beats, double
  tempo~\cite{ahn2026smc} \reported.
\item Solo-guitar recordings without accompaniment: about half the beats
  missed (Table~\ref{tab:bench-beats}).
\end{itemize}

% !TeX root = ../main.tex
\section{String and fret assignment (fingering)}
\label{sec:fingering}

\subsection{What the stage must do}

Input: notes (onset, pitch, duration), grouped into chords, plus a tuning.
Output: one string and fret per note, such that every chord is playable (one
note per string, a reachable fret span) and the sequence is comfortable and
idiomatic. Section~\ref{sec:problem} showed why this is ambiguous: a pitch has
up to five positions, a three-note chord can have dozens of shapes.

\subsection{Search with hand-designed costs}

\begin{keyidea}[: fingering as a shortest path]
Treat every possible position (or chord shape) of every note as a node, give
each node a cost (high frets, stretched shapes) and each move between
consecutive nodes a cost (hand movement), and find the cheapest path through
the whole piece with the Viterbi algorithm (Section~\ref{sec:viterbi}). Sayegh
framed it this way in 1989~\cite{sayegh1989fingering}; HMM versions decode
left-hand forms~\cite{hori2016minimax}.
\end{keyidea}

The weakness is that the costs are local and hand-made. They know nothing of
long-range habits: a riff repeated in another position, idiomatic chord shapes,
an open string used as a pedal tone, or the way solos move up the neck.

\subsection{Learned fingering}

\paragraph{Symbolic models} learn from human tabs (no audio). They treat
fingering as translation from note tokens to string/fret tokens, usually
pre-trained on DadaGP. MIDI-to-Tab~\cite{edwards2024miditotab} (encoder-decoder
transformer, masked pre-training, fine-tuned on professional transcriptions)
was rated more playable than other algorithms in a user study \reported.
Fretting-Transformer~\cite{hamberger2025fretting} (T5, with tuning and capo
conditioning) beat an A* baseline and Guitar Pro's own conversion \reported.
Hsu et al.~\cite{hsu2026noteevent} added note-event tokens and \term{constrained
decoding} that forbids string-fret pairs inconsistent with the pitch, and beat
Fretting-Transformer, especially with little training data \reported. None of
these three has released weights.

\paragraph{Audio-conditioned models} also listen: the same pitch sounds
different on a thick and a thin string. TART's third stage is a T5
conditioned on audio features~\cite{gupta2026tart}; with its note stage it
reaches 71.8\% string-fret Tab F1 averaged over its benchmarks, 8.5 points above
the best earlier baseline \reported. No code was found.

\begin{inapp}[: the Viterbi (fallback) \done]
\file{transcribe/arrange.rs}, \code{fret_events}. States are fingerings of each
chord (or single note); the hand position is carried through open strings. The
cost of a fingering is a weighted sum (\code{FingeringCosts}, weights fitted
with \code{eval_fingering} on GuitarSet players 00--02 and 10 human Guitar Pro
tabs):
\begin{center}
\small
\begin{tabular}{@{}lrl@{}}
\toprule
Term & Weight & Per \\
\midrule
\code{span} & 0.40 & fret between lowest and highest fretted note of a chord \\
\code{position} & $-0.01$ & fret of hand position (almost neutral after fitting) \\
\code{high} & 0.08 & fret above the 12th \\
\code{open} & 0.06 & open string \\
\code{gap} & 0.80 & unused string inside a chord (must be muted) \\
\code{moving} & 0.28 & fret of hand movement between events \\
\code{prior} & 0.11 & $-\ln P(\text{string}\mid\text{pitch})$, learned from 34\,150 GuitarSet notes \\
\bottomrule
\end{tabular}
\end{center}
Two hard-ish rules come from the techniques stage: a slide or hammer-on that
changes string costs 4.0 (\code{LINK_STRING_CHANGE}), a bend on an open string
6.0 (\code{BENT_OPEN_STRING}). Before fingering, a tuning is chosen from the
lowest notes (standard, E$\flat$, drop D, drop C, 7 and 8 strings).
\end{inapp}

\begin{inapp}[: the learned tagger \done{} (local model)]
\label{sec:fingering-app}
\file{transcribe/fingering_net.rs}, trained by \file{scripts/train/fingering/}.
\begin{itemize}
\item \textbf{Model}: a 4-layer Transformer \emph{tagger} (0.57\,M parameters).
  It does not generate tokens; for each note it outputs a score for each of
  the 6 strings. Inputs per note: the fret the note would have on each string
  (or ``impossible''), plus 16 timing and interval features (\code{common.features}).
  Windows of 128 notes.
\item \textbf{Decoder}: puts each chord on distinct strings, adds a stretch
  cost, and runs a hand-position Viterbi over the tagger's scores, with the
  same technique rules as above.
\item \textbf{Training data} (``ship'' variant, about 70\,k notes): GuitarSet
  00--02, EGDB train, Guitar-TECHS P1/P2. Settings were chosen with GuitarSet
  player 02 held out; never on 03--05.
\item \textbf{Use}: 6-string guitars, when \code{fingering_tagger.onnx} is in the
  models folder and the Transcribe form (or \cmd{--fingering learned}) asks for
  it; otherwise, and on any model error, the Viterbi. Local only (trained on
  EGDB).
\item \textbf{Parity}: the Rust port agrees with the Python model on 99.94\% of
  41\,293 test notes \measured. It needs each note's own onset: giving every
  chord note the chord's onset costs 0.02 accuracy.
\item \textbf{Cost}: about 6--10\ms{} per 128-note window on the CPU.
\end{itemize}
\end{inapp}

\subsection{What we measured}

\begin{table}[h]
\centering
\caption{String accuracy given the true notes, all test items, mean of 3 seeds
for the tagger \measured{} (\code{scripts/train/fingering/}). ``Full'' adds GOAT
training data (non-commercial).}
\label{tab:fingering}
\small
\begin{tabular}{@{}lccccc@{}}
\toprule
Method & GuitarSet comp & GuitarSet solo & EGDB test & EGSet12 & GOAT test \\
\midrule
Lowest fret (naive) & 0.33 & 0.25 & 0.65 & 0.43 & -- \\
Viterbi (app fallback) & 0.807 & 0.549 & 0.677 & 0.688 & 0.648 \\
Tagger, ship (app) & \textbf{0.853} & \textbf{0.616} & \textbf{0.806} & 0.698 & 0.667 \\
Tagger, full & 0.845 & 0.613 & 0.783 & \textbf{0.722} & \textbf{0.674} \\
\bottomrule
\end{tabular}
\end{table}

With transcribed notes (v2 CRNN) the tagger raises Tab F1 from 0.53--0.60 to
0.64--0.68 on the EGDB tones, from 0.53 to 0.63 on GuitarSet solo, and from
0.59 to 0.65 on BS-RoFormer stems of band mixes \measured. In the app pipeline
(\code{eval_guitarset}, players 03--05), the median share of notes on the
reference string and fret goes from 68\% to 78\% with GAPS notes on clean
guitar, and from 67\% to 76\% with v2 on ``Metal High Gain'' \measured.
Chords stay playable: fewer than 0.5\% impossible stretches, except 6.5\% on
EGSet12.

\subsection{Failure modes}
\begin{itemize}
\item \textbf{Too few players.} The tagger scores 0.95 and more on its training
  players and 0.57--0.69 on new ones; a 3.3\,M-parameter version did no better.
  Player diversity, not model size, is the limit \measured. DadaGP
  pre-training is the planned fix.
\item \textbf{Transposition hurts.} Training on passages shifted by $k$ frets
  lowered solo accuracy to 0.43: positions depend on the key \measured.
\item \textbf{Solos.} The Viterbi puts only 55--56\% of solo notes on the
  player's string (80\% for accompaniment): it keeps notes low while solos move
  up the neck.
\item \textbf{Licence.} The tagger is trained on EGDB, so it cannot be
  distributed until EGDB's authors allow it or it is retrained without EGDB.
\end{itemize}

% !TeX root = ../main.tex
\section{Playing techniques}
\label{sec:techniques}

\subsection{What the stage must do}

Mark, per note, how it is played: bend (and how far, in quarter tones) and
release, vibrato, slide into or out of the note, hammer-on and pull-off (a note
sounded by the fretting hand without picking), palm mute, natural and
artificial harmonics, dead (muted) notes. Some of these change the fingering:
a slide or a hammer-on stays on one string, an open string cannot be bent.

\subsection{How techniques are detected}

\begin{keyidea}[: read the pitch contour]
Most techniques are visible in the fine pitch contour of a single note
(Section~\ref{sec:audio}): a bend is a smooth rise and maybe a fall back, vibrato
a periodic wobble, a slide a run of semitone steps, a hammer-on a pitch jump
\emph{without} a new pick attack. Rules on the contour and the attack can
therefore detect them on clean single-note lines.
\end{keyidea}

Chen et al.~\cite{chen2015technique} introduced rule-based detection of bends,
vibrato, slides and hammer-ons on real electric guitar recordings.
TENT~\cite{su2019tent} built technique recognition into note tracking for
monophonic solos, with a 14.7\% higher note F-score than an earlier method
\reported. TART~\cite{gupta2026tart} classifies techniques per note with a small
network (an MLP). GOAT~\cite{loth2025goat} and IDMT-SMT-Guitar~\cite{kehling2014idmt}
provide labels. Almost all evaluation is on clean, monophonic material: distortion
blurs the contour and compression hides the pick attack, so techniques on
amplified rhythm parts are largely untested.

\begin{inapp}[: rules on Basic Pitch's contour \done]
\file{transcribe/techniques.rs}; \code{GVA_TECHNIQUES=0} turns it off.
\begin{itemize}
\item \textbf{Contour}: Basic Pitch's third output, 264 bins at three per
  semitone (MIDI pitch $p$ peaks at bin $3(p-21)+1$, measured). With the CRNN
  engine, Basic Pitch is run again on the guitar audio just for this contour.
\item \textbf{Single-note passages only}; chords are skipped.
\item \textbf{Split bends merged}: the note model often reports a bend as two
  notes; a glide of at least 30\ms{} that does not step through semitones joins
  them back into one bent note.
\item \textbf{Bend and release} inside a note; \textbf{vibrato} at 3.5--9\Hz,
  at least 0.25 semitone deep, at least 2 cycles.
\item \textbf{Slides}: stepped glides, downward glides or glides of more than
  3 semitones, with a gap of up to 120\ms{} between the notes.
\item \textbf{Hammer-on / pull-off}: a pitch jump within 40\ms{} with no
  high-frequency pick burst (\code{PICK_ATTACK}).
\item \textbf{Output}: \code{arrange} keeps techniques per pitch; the tab gets
  \code{bend}, \code{bend_value} (quarter tones) and \code{bend_release} on the
  struck note, \code{slide} and \code{hammer} on the last tied piece, vibrato
  on all pieces. The GP5 export writes bends and shift slides, and the importer
  reads bends back.
\item Thresholds can be swept with \code{GVA_BEND_RISE}, \code{GVA_BEND_START},
  \code{GVA_BEND_GLIDE}.
\end{itemize}
\end{inapp}

\subsection{What we measured}

\begin{itemize}
\item On GuitarSet the rules change nothing measurable (F1 0.83 with or without
  them; 0.76 either way on the crunch render): the 60 excerpts contain only 19
  annotated glides \measured.
\item On real technique labels (IDMT, Guitar-TECHS, GOAT, via
  \code{eval_manifest} with a 5th column in the notes file) \measured:
  vibrato is found on about 20--40\% of labelled vibrato notes, bends on
  35--56\%, slides on about 30\% (often marked as legato instead), and false
  alarms on plain notes are 0--4\%.
\end{itemize}

So the rules are cautious: they rarely invent a technique, and they miss most.

\subsection{Failure modes}
\begin{itemize}
\item No palm mute, harmonics, dead notes or tapping (not implemented).
\item Chords are skipped by design.
\item Under distortion the contour gets noisy and compressed attacks hide the
  pick burst, so bends, vibrato and hammer-ons are missed exactly where they are
  most common \expected.
\end{itemize}

\begin{nextsteps}
Tune the hammer-on attack threshold and the bend rules on GOAT and
Guitar-TECHS; then a learned per-note classifier (task T3 in
Section~\ref{sec:next}) for palm mute, harmonics and dead notes, keeping the
rules for any class where the model is worse.
\end{nextsteps}

% !TeX root = ../main.tex
\section{End-to-end audio-to-tablature models}
\label{sec:e2e}

\subsection{The idea}

An end-to-end model predicts string and fret directly from audio, so the
fingering can use timbre (which string a note sounds like) and the constraints
of the fretboard can shape the transcription itself.

\subsection{The main models}

\begin{description}[style=nextline,leftmargin=1.2cm]
\item[Before deep learning] Barbancho et al.~\cite{barbancho2012chords}
  transcribed chords and their fingering with an HMM over fingering
  configurations.
\item[TabCNN (2019)] \cite{wiggins2019tabcnn} A small CNN on a constant-Q
  spectrogram that outputs, for each frame and each of the six strings, a
  probability over ``silent, open, fret 1, \ldots, fret 19''. No note-level
  decoding and no temporal model. Trained and tested on GuitarSet; the standard
  end-to-end baseline. Cwitkowitz et al.~\cite{cwitkowitz2022inhibition} added a
  loss that discourages note pairs unlikely on the fretboard.
\item[FretNet (2023)] \cite{cwitkowitz2023fretnet} Predicts continuous pitch
  contours grouped by string and fret, so bends and vibrato can be represented
  (a semitone grid cannot). Code but no weights released.
\item[SynthTab (2024)] \cite{zang2024synthtab} Pre-trains on DadaGP tabs
  rendered through guitar plugins; fine-tuning then gives better results and
  less overfitting across datasets. The main evidence that data scale, not
  architecture, limits end-to-end models.
\item[GOAT (2025)] \cite{loth2025goat} Preliminary tab-transcription results on
  its electric DI and amplifier data; the first large source of electric guitar
  paired with technique-annotated tabs.
\item[Noise2Fret (2026)] \cite{simionato2026noise2fret} A diffusion model over a
  continuous representation of string and fret, conditioned on spectral
  features, with five auxiliary losses for playability (pitch-class distance,
  positional distance, circle-of-fifths distance, string similarity, hand-span
  feasibility). Beats its baselines on GuitarSet and GOAT \reported.
\item[TART (2025--2026)] \cite{gupta2025tart,gupta2026tart} Presented as
  end-to-end but actually four stages: audio to MIDI, a technique classifier,
  audio-conditioned fingering, and a rule-based tab generator. Zero-shot on
  GuitarSet, EGDB and noisy versions of both, it reports 81.35\% note F50,
  71.8\% string-fret Tab F1 and 54.08\% end-to-end Tab F1 on average
  \reported. The drop from 71.8 to 54.08 shows how errors compound across
  stages.
\end{description}

All of these take an isolated guitar as input. None is trained or evaluated on
guitar inside a band mix, and most are trained on GuitarSet, which is acoustic.

\begin{inapp}[: not used, but benchmarked]
The app has no end-to-end model. We ran the two TabCNN variants with released
weights (TabCNN on GuitarSet, and TabCNN-GPFX, its tone-augmented variant from
Pedroza et al.~\cite{pedroza2024tones}) against our modular pipeline; the
results are in Section~\ref{sec:bench-tab}. TabCNN wins on GuitarSet solo lines
and loses everywhere on electric guitar and mixes.
\end{inapp}

\subsection{Noise2Fret under its own protocol}

We ran Noise2Fret's released checkpoints (CC BY 4.0) with its own code and
protocol (GuitarSet player 00 held out, 100\ms{} windows): pitch F 0.924 / 0.926
and Tab F 0.921 / 0.922 without / with the auxiliary losses \measured.

\begin{caution}[: why 0.92 is not comparable]
That protocol scores only 100\ms{} windows that \emph{contain a note onset}
(\code{drop_empty=True}), so silent passages can never produce false notes, and
timing is judged per window instead of with a 50\ms{} onset tolerance. The number
cannot be compared with the note-level Tab F1 in Section~\ref{sec:bench-tab}.
Running the model on whole recordings would need an onset detector it does not
include; its event-based follow-up targets that.
\end{caution}


\part{Guitar Virtual Amp: system and measurements}
% !TeX root = ../main.tex
\section{The app's pipeline end to end}
\label{sec:app}

This section puts the stages of Part~II together as they run in the app.
Figure~\ref{fig:app-pipeline} is the data flow; Table~\ref{tab:app-components}
lists each component, its file and what happens when an optional part is
missing.

\begin{figure}[h]
\centering
\begin{tikzpicture}[
  node distance=3.5mm,
  box/.style={draw, rounded corners=2pt, font=\small, align=center,
              minimum height=8mm, text width=40mm, inner sep=3pt, fill=white},
  opt/.style={box, dashed},
  arr/.style={-{Stealth[length=2mm]}, semithick}]
  \node[box] (in) {Audio file or YouTube\\(\code{yt-dlp})};
  \node[box, below=of in] (dec) {Decode + SHA-256};
  \node[opt, below=of dec] (sep) {Separation\\BS-RoFormer / Demucs};
  \node[opt, right=12mm of sep] (beat) {Beats, bars, meter\\Beat This! / DSP};
  \node[box, below=of sep] (pan) {Pan split\\(left / centre / right)};
  \node[box, below=of pan] (notes) {Guitar notes per part\\GAPS / v2 / Basic Pitch};
  \node[box, below=of notes] (tech) {Techniques\\(Basic Pitch contour)};
  \node[box, left=12mm of notes] (bass) {Bass notes\\Basic Pitch};
  \node[box, below=of tech] (arr) {Arrange: quantise,\\fingering, tuning, notation};
  \node[box, below=of arr] (out) {GP5 tab + \code{.gvtab.json}\\in the tab repository};
  \node[opt, right=12mm of out] (tone) {Tone estimate per guitar\\(in progress, optional)};
  \draw[arr] (in) -- (dec);
  \draw[arr] (dec) -- (sep);
  \draw[arr] (dec.east) -| (beat.north);
  \draw[arr] (sep) -- (pan);
  \draw[arr] (pan) -- (notes);
  \draw[arr] (notes) -- (tech);
  \draw[arr] (tech) -- (arr);
  \draw[arr] (sep.west) -| (bass.north);
  \draw[arr] (bass.south) |- (arr.west);
  \draw[arr] (beat.south) |- (arr.east);
  \draw[arr] (arr) -- (out);
  \draw[arr] (out) -- (tone);
\end{tikzpicture}
\caption{Data flow of \code{transcribe::run} (\file{transcribe/mod.rs}). Dashed
boxes are optional: the pipeline still runs without them, with the fallbacks of
Table~\ref{tab:app-components}.}
\label{fig:app-pipeline}
\end{figure}

\begin{table}[h]
\centering
\caption{Components of the transcription pipeline.}
\label{tab:app-components}
\small
\begin{tabularx}{\textwidth}{@{}>{\raggedright}p{2.2cm}>{\raggedright}p{3.6cm}>{\raggedright}p{3.6cm}X@{}}
\toprule
Stage & Default & File & Without the optional part \\
\midrule
Source & file (mp3, flac, ogg, wav, m4a) or YouTube & \file{transcribe/youtube.rs} & -- \\
Separation & BS-RoFormer SW (MSST, Python) & \file{transcribe/separation.rs} & Demucs \code{htdemucs_6s}; else the full mix, one guitar track \\
Guitar parts & pan split & \file{transcribe/stereo.rs} & one guitar track (mono input or no separation) \\
Beats and bars & Beat This! final0 (small0) & \file{transcribe/beat_net.rs} & DSP tracker, 4/4 only (\file{beats.rs}) \\
Guitar notes & GAPS CRNN; v2 if chosen & \file{transcribe/guitar_crnn.rs} & Basic Pitch \\
Bass notes & Basic Pitch & \file{transcribe/basic_pitch.rs} & -- \\
Techniques & contour rules & \file{transcribe/techniques.rs} & -- \\
Fingering & learned tagger (6 strings) & \file{transcribe/fingering_net.rs} & Viterbi (\file{arrange.rs}) \\
Notation & quantise, tuplets, ties, tuning & \file{transcribe/arrange.rs} & -- \\
Output & GP5 + sidecar & \file{shared/tab/export.rs}, \file{repository.rs} & -- \\
\bottomrule
\end{tabularx}
\end{table}

\subsection{Entry points}

\begin{itemize}
\item \textbf{Tabs page}, button ``Transcrever áudio'', and the Spotify page
  for a track without a tab (\file{frontend/ui/transcribe.rs}). The form lets
  the user pick a file or search YouTube, and choose the guitar model and the
  fingering engine.
\item \textbf{Batch CLI}, which skips audio already in the repository (also
  \code{--youtube}):
\begin{verbatim}
cargo run --release -p guitar_desktop --bin transcribe -- <paths>
\end{verbatim}
\item \textbf{YouTube} (\file{transcribe/youtube.rs}): wraps the external
  \code{yt-dlp}; searches without an API key, ranks results by title words,
  ``Artist -- Topic'' channels, penalties for live/cover/karaoke and closeness
  to the Spotify duration; downloads audio into a private cache, never into
  the repository. Video IDs are validated before reaching the command line.
\end{itemize}

\subsection{What a transcription produces}

The result goes to the local tab repository
(\file{<data dir>/guitar_virtual_amp/tab_repository/}):
\begin{itemize}
\item \file{Artist - Title.gp5}, a Guitar Pro 5 file written through the
  \code{guitarpro} crate (opens in Guitar Pro, TuxGuitar, MuseScore). Tempos are
  rounded per bar, as GP5 requires.
\item \file{Artist - Title.gvtab.json}, the \term{sidecar} (\code{TabMetadata}):
  the source audio's SHA-256 (to skip or replace repeats), the engines used
  (separator, note model, beat model, fingering), the YouTube ID when there is
  one, and the time in the audio of every bar (\code{bar_audio_secs}). The tab
  player uses the bar times (\code{TabMetadata::sync_map}) to play the original
  recording in exact sync, whatever the GP5 tempo rounding.
\end{itemize}

\subsection{Cost}

On the laptop (RTX 4070, 8\,GB): a 4.4-minute song takes about 70\,s and
2.4\,GB of RAM \measured; BS-RoFormer separation alone is 34\,s on a 261\,s
song. The CRNN windows are the memory peak (about 600\,MB each in tract, so at
most 4 at a time). The tract crates are built with \cmd{opt-level = 3} even in
debug builds, otherwise transcription takes minutes.

\begin{inapp}[: tone estimate, experimental \done]
An optional last stage: with
``Estimar o timbre (experimental)'' checked in the Transcribe form (setting
\code{AppSettings::transcription_tone}), each guitar part's audio and notes go
to \code{tone::estimate_from_transcription}, and the estimated chain is saved as
a user preset tagged with the song and the track (\code{tone::save_preset}), so
the Tabs page can suggest it. It needs separation. The method is the subject of
the companion document \file{docs/papers/tone-estimation/}.
\end{inapp}

% !TeX root = ../main.tex
\section{Our benchmarks}
\label{sec:benchmark}

Published numbers use different test sets and metrics
(Section~\ref{sec:eval}), so the project runs its own benchmarks. Stage-level
results were already given in Part~II; this section explains the test sets and
protocols and gives the results that involve several stages.

\subsection{Test sets and what each one tests}

\begin{table}[h]
\centering
\caption{The benchmark matrix. Training sets keep their test parts out.}
\label{tab:testsets}
\small
\begin{tabularx}{\textwidth}{@{}>{\raggedright}p{3.4cm}>{\raggedright}p{4.2cm}>{\raggedright}p{2.6cm}X@{}}
\toprule
Set & Audio & Labels & Tests \\
\midrule
GuitarSet players 03--05 & solo acoustic, mic & notes + strings & clean baseline, fingering \\
GuitarSet $\times$ presets (\code{GVA_EVAL_TONE}) & same, through the app's chains & same & robustness to tone \\
EGDB test & electric DI, Marshall, Mesa, Plexi sims, 5 real clips & notes + strings & amp tones, fingering on electric \\
EGSet12 & real electric into a miked amp (12) & notes + strings & real electric \\
Guitar-TECHS P3 + techniques & electric DI and miked amp & MIDI notes, techniques & real electric, techniques \\
IDMT-SMT-Guitar & DI notes and licks & notes + style & techniques \\
GOAT test & DI + 5 amp renders (7 items per tone) & GP tabs with techniques & Tab F1, techniques \\
Band mixes (\code{make_mixes.py}) & labelled guitar over MUSDB18 drums, bass, vocals & notes & separation, full mix \\
MoisesDB (\code{prepare_moisesdb.py}) & 222 real songs, 60\,s each & none (isolated guitars) & separation on real productions \\
GP SoundFont renders (\code{eval_transcription}) & the user's Guitar Pro files, synthesised & tabs & whole pipeline, several tracks \\
Real songs (\code{align_tabs.py}) & commercial recordings & the user's tabs, aligned & the actual use case (not collected yet) \\
\bottomrule
\end{tabularx}
\end{table}

\subsection{Protocols}

\paragraph{Stage benchmarks (Python).} \code{scripts/train/bench_*.py} run one
model per stage on up to 12 items per group (spread evenly; GOAT has 7 per tone,
EGDB real 5), on a DGX Spark (NVIDIA GB10 GPU). Metrics come from
\code{mir_eval}, means per item. Guitar-TECHS and GOAT labels are shifted per
recording to remove a constant offset. Fingering is scored by running the app's
own Viterbi or tagger through \cmd{cargo run --example fingering}, so the Python
and Rust numbers use the same code.

\paragraph{End-to-end benchmarks (Rust).} Ignored tests in
\file{transcribe/eval.rs} run \code{transcribe::run} (the real app pipeline) and
score the finished tab: \code{eval_guitarset}, \code{eval_transcription},
\code{eval_fingering}, \code{eval_real}, \code{eval_guitar_pairs},
\code{eval_manifest}. Commands are in Appendix~\ref{app:code}.

\begin{caution}[: known biases of these benchmarks]
\begin{itemize}
\item The band mixes put a guitar part over an \emph{unrelated} song (different
  key and tempo). Models that use musical context (sequence-to-sequence
  transcribers) may be penalised.
\item Both TabCNN models were trained on GuitarSet, probably including players
  03--05, so their GuitarSet numbers may be optimistic.
\item MuScriptor's training set is private; overlap with our test sets cannot be
  ruled out.
\item v2 saw other clips and players of EGDB, Guitar-TECHS and GuitarSet (never
  GOAT, EGSet12 or MUSDB18 audio). EGDB DI/Plexi/real was its validation set,
  so its EGDB scores are slightly optimistic.
\item The Viterbi's weights were fitted on GuitarSet 00--02 and 10 human tabs,
  and its string prior comes from GuitarSet 00--02, so string agreement on
  those is optimistic.
\end{itemize}
\end{caution}

\subsection{Tablature: modular against end-to-end}
\label{sec:bench-tab}

\begin{table}[h]
\centering
\caption{Tab F1 (F50 note match on the same string) \measured{}
(\code{bench_tab.py}). Modular rows place the transcribed notes with the
Viterbi. The last row gives the Viterbi's string accuracy on the true notes.
EGDB amps: mean of three amp sims; mix columns: EGDB clips in band mixes.
$^{*}$String accuracy, not Tab F1.}
\label{tab:bench-tab}
\footnotesize
\setlength{\tabcolsep}{3pt}
\begin{tabular}{@{}lccccccc@{}}
\toprule
Method & EGDB DI & EGDB amps & EGSet12 & GS comp & GS solo & Raw mix & BS-RoF.\ stem \\
\midrule
GAPS + Viterbi & 0.52 & 0.44 & 0.57 & 0.71 & 0.55 & 0.36 & 0.47 \\
v2 + Viterbi & \textbf{0.60} & \textbf{0.55} & 0.57 & 0.73 & 0.53 & 0.33 & \textbf{0.59} \\
Basic Pitch + Viterbi & 0.45 & 0.43 & 0.40 & 0.46 & 0.42 & 0.28 & 0.45 \\
YourMT3+ + Viterbi & 0.59 & 0.48 & 0.30 & \textbf{0.76} & 0.54 & \textbf{0.43} & 0.46 \\
MuScriptor large + Viterbi & 0.53 & 0.53 & \textbf{0.63} & 0.64 & 0.50 & 0.37 & 0.54 \\
\midrule
TabCNN (end-to-end) & 0.40 & 0.32 & 0.31 & 0.71 & 0.76 & 0.18 & 0.34 \\
TabCNN-GPFX (end-to-end) & 0.43 & 0.37 & 0.40 & 0.74 & \textbf{0.81} & 0.20 & 0.37 \\
\midrule
Viterbi on true notes$^{*}$ & \multicolumn{2}{c}{0.74} & 0.69 & 0.80 & 0.56 & -- & -- \\
\bottomrule
\end{tabular}
\end{table}

Reading Table~\ref{tab:bench-tab}:
\begin{itemize}
\item \textbf{The modular pipeline wins on electric guitar and mixes} (EGDB amps
  0.55 against 0.37; separated mix 0.59 against 0.37). TabCNN's weak point is
  note detection on electric guitar (its F50 is 0.36--0.50 on EGDB), not string
  choice.
\item \textbf{TabCNN wins on GuitarSet solo lines} (0.81 against 0.53). The
  cause is the Viterbi: given the true notes it puts only 56\% of solo notes on
  the player's string. This was the measurement that made learned fingering the
  priority; with the tagger, v2's Tab F1 on GuitarSet solo rises to 0.63 and on
  BS-RoFormer stems to 0.65 (Section~\ref{sec:fingering}).
\item \textbf{MuScriptor large gives the best tab on EGSet12} (0.63 against 0.57
  at equal F50): fewer stray notes inside chords, so the fingering stage gets
  cleaner chord shapes.
\end{itemize}

\subsection{Real electric guitar: EGSet12}

\begin{table}[h]
\centering
\caption{EGSet12, 12 real solo electric performances (Telecaster into a miked
amp, 380\,s), none used in training \measured.}
\label{tab:egset12}
\small
\begin{tabular}{@{}lccc@{}}
\toprule
Method & F50 & Onset-offset F1 & Tab F1 \\
\midrule
GAPS (+ Viterbi) & \textbf{0.89} & 0.41 & 0.57 \\
v2 (+ Viterbi) & 0.88 & \textbf{0.57} & 0.57 \\
MuScriptor large (+ Viterbi) & \textbf{0.89} & 0.56 & \textbf{0.63} \\
Basic Pitch (+ Viterbi) & 0.76 & 0.56 & 0.40 \\
YourMT3+, guitar classes (+ Viterbi) & 0.54 & 0.45 & 0.30 \\
TabCNN / TabCNN-GPFX & 0.48 / 0.56 & -- & 0.31 / 0.40 \\
\bottomrule
\end{tabular}
\end{table}

TabCNN-GPFX's paper reports 0.585 \emph{frame-level} tab F1 on EGSet12
\reported; at note level it reaches 0.40 here, a reminder that frame and note
metrics differ.

\subsection{Real mixes: MoisesDB}
\label{sec:bench-moises}

MoisesDB songs come with isolated guitar tracks but no note labels. The
benchmark (\code{bench_moisesdb.py}) cuts the 60\,s with the most guitar from
each of the 222 songs with guitar (113 with distorted guitar, 66 clean
electric, 43 acoustic only), transcribes the mix or a separated stem, and scores
the notes against notes transcribed from the isolated guitars.

\begin{keyidea}[: a benchmark without labels]
If a model's notes on the separated stem agree with its own notes on the clean
isolated guitar, the mix and the separator cost it little. Two references are
used: \term{self} (the same model's notes on the isolated guitars: what the mix
costs that model) and \term{consensus} (notes on which several models agree).
This ranks \emph{separators} and mix robustness well. It does not rank note
models, because each model scores best against a reference that includes
itself.
\end{keyidea}

\begin{table}[h]
\centering
\caption{MoisesDB, all 222 songs, F50 against the consensus of v2 and GAPS on
the isolated guitars (``self'' reference in brackets) \measured.}
\label{tab:moises}
\small
\begin{tabular}{@{}lcc@{}}
\toprule
Input & GAPS & v2 \\
\midrule
Raw mix & 0.42 (0.45) & 0.41 (0.45) \\
\code{htdemucs_6s} stem & 0.58 (0.63) & 0.58 (0.65) \\
BS-RoFormer SW stem & \textbf{0.71 (0.76)} & \textbf{0.70 (0.78)} \\
\bottomrule
\end{tabular}
\end{table}

On the 24 songs MuScriptor large ran on (it is too slow for all 222), it loses
less on the raw mix than the CRNNs (self reference 0.54 against 0.43--0.45) but
still scores only 0.42 there, against 0.71 for v2 on the BS-RoFormer stem; on
the stem it is no better than the CRNNs \measured. The question left open by the
synthetic mixes is settled: keep separator + CRNN. By guitar kind, v2 on the
BS-RoFormer stem scores 0.56 on distorted, 0.75 on clean electric and 0.83 on
acoustic guitar.

\subsection{The app end to end}
\label{sec:bench-app}

\begin{itemize}
\item \textbf{GuitarSet players 03--05} (\code{eval_guitarset}, final-tab F1
  median): GAPS 0.84 clean, 0.76 ``Blues Crunch'', 0.43 ``Metal High Gain''; v2
  0.86, 0.81, 0.69 (Section~\ref{sec:notes}). String/fret agreement with the
  tagger: 78\% (GAPS, clean), 76\% (v2, metal).
\item \textbf{History} on an earlier 60-excerpt GuitarSet selection (\file{ROADMAP.md}):
  final-tab F1 0.70 with the first pipeline, 0.73 with Basic Pitch and the new
  rhythm stage, 0.83 with the guitar CRNN. Guitar Pro SoundFont renders of
  isolated guitar: 0.82 $\to$ 0.85; full mixes without separation: 0.52 $\to$
  0.68.
\item \textbf{Two guitars in one stereo mix} (\code{eval_guitar_pairs}): lead
  part 0.85--0.88 with the pan split against 0.40--0.43 as one merged track
  (Section~\ref{sec:separation}).
\item \textbf{Viterbi fingering, first cost function $\to$ fitted costs and
  string prior} (true notes, share on the human's string): GuitarSet players
  03--05 65.9\% $\to$ 75.4\%; 10 human Guitar Pro tabs 58.5\% $\to$ 65.4\%
  (optimistic: the weights were fitted on those tabs).
\end{itemize}

\subsection{Toward a real-song benchmark}
\label{sec:bench-real}

The goal is 30--50 commercial recordings the owner has, each with a Guitar Pro
tab of the same song, aligned to the audio and hand-checked.
\code{scripts/train/align_tabs.py} \done{} aligns a tab to a recording with
\code{synctoolbox} (multi-resolution DTW on chroma and onset features computed
from the tab's notes; the best transposition within $\pm2$ semitones is
applied), writes per-track notes in recording time, and a hand-check copy of the
tab with a sidecar so the Tabs page plays the recording in sync.
\code{guitar_core/examples/tab_sync_check.rs} checks such a copy through the
app's own loader. Only pairs marked \code{checked} go into \code{eval_real.tsv}.

Validation on recordings with known note times \measured{}
(\code{validate_tab_align.py}): note F1 at 50\ms{} between the aligned tab and
the labels is 0.92 on EGSet12 (median onset error 12\ms) and 0.80 on the GOAT
test split (18\ms); with per-bar tempo drift plus an intro added, 0.84 and 0.73.
The tab's own timing on these click-tracked sets gives 0.95 and 0.92, the
ceiling. Rejection rules (\code{cost_ratio} $>0.70$ or \code{stall} $>0.10$)
kept 121 of 128 correct pairs and rejected all 168 wrong-tab pairs and all 28
tabs with 15--25\% of their bars removed, with a thin margin.

\begin{caution}[: status]
No real song pairs have been collected or checked yet, and \code{eval_real} is
written but unmeasured. Until then, every number in this document is on solo
guitar, synthetic mixes or label-free real mixes. The use case itself is not
yet measured.
\end{caution}

% !TeX root = ../main.tex
\section{Assessment: where we stand and what comes next}
\label{sec:assessment}

\subsection{The app against the state of the art}

The literature splits in two (Table~\ref{tab:methods}). Methods that handle real
mixes (MuScriptor, YourMT3+, source separators) output notes or stems without
fingering; methods that output tablature (TabCNN, FretNet, Noise2Fret, TART)
require a solo guitar. A modular pipeline that chains a separator, a guitar note
model and a fingering model is currently the only way to get a tab from a
commercial recording. The app is such a pipeline, built from the strongest
components with released weights, and our measurements support the choice: on
electric guitar and on separated mixes it beats the end-to-end models we could
run by 0.17--0.23 Tab F1 (Table~\ref{tab:bench-tab}).

The most direct published comparison is TART: 54.08\% end-to-end Tab F1 averaged
over clean and noisy solo-guitar sets \reported. Our closest numbers are Tab F1
0.64--0.68 on the EGDB tones and 0.65 on separated band mixes with v2 and the
tagger \measured. These are \emph{not} the same test sets, so the only safe
conclusion is that the app is in the same range as the best published system,
on harder inputs (amplified and mixed guitar).

\begin{table}[h]
\centering
\caption{Main methods by stage. ``Mix'': designed or shown to work on
multi-instrument recordings. ``Here'': whether we ran it. Reported results use
different test sets.}
\label{tab:methods}
\small
\setlength{\tabcolsep}{3pt}
\begin{tabularx}{\textwidth}{@{}>{\raggedright}p{3.0cm}>{\raggedright}p{1.7cm}>{\raggedright}p{2.6cm}c>{\raggedright}p{2.6cm}X@{}}
\toprule
Method & Stage & Output & Mix & Here & Reported (as published) \\
\midrule
HT Demucs~\cite{rouard2023htdemucs} & separation & 4 or 6 stems & yes & app fallback & 9.20\dB{} SDR, MUSDB18 \\
BS-RoFormer~\cite{lu2023bsroformer} & separation & stems & yes & app default & 9.80\dB{} SDR; SDX23 winner \\
Banquet~\cite{watcharasupat2024banquet} & separation & queried stem & yes & not yet & beats \code{htdemucs_6s} on guitar \\
GAPS model~\cite{riley2024gaps} & notes & notes & no & app default & SOTA on GuitarSet \\
Basic Pitch~\cite{bittner2022basicpitch} & notes & notes, contour & no & bundled & lightweight baseline \\
YourMT3+~\cite{chang2024yourmt3} & notes & notes + instrument & yes & benchmarked & limited on pop \\
MuScriptor~\cite{rouard2026muscriptor} & notes & notes + instrument & yes & benchmarked & 60.4 onset F1 (own test) \\
Beat This!~\cite{foscarin2024beatthis} & rhythm & beats, downbeats & yes & app default & SOTA beat F1 without DBN \\
AlignBeat~\cite{ahn2025alignbeat} & rhythm & beats, downbeats & yes & not yet & +3.6 downbeat CMLt \\
Viterbi / HMM~\cite{sayegh1989fingering,hori2016minimax} & fingering & string, fret & -- & app fallback & hand-made costs \\
MIDI-to-Tab~\cite{edwards2024miditotab} & fingering & string & -- & no weights & preferred in a user study \\
Fretting-Transformer~\cite{hamberger2025fretting} & fingering & string, fret & -- & no weights & beats A* and Guitar Pro \\
Chen et al.~\cite{chen2015technique}, TENT~\cite{su2019tent} & techniques & per note & no & rules like these & TENT +14.7\% note F \\
TabCNN~\cite{wiggins2019tabcnn} & end-to-end & frets per frame & no & benchmarked & GuitarSet baseline \\
Noise2Fret~\cite{simionato2026noise2fret} & end-to-end & string, fret & no & own protocol only & beats baselines \\
TART~\cite{gupta2026tart} & hybrid & tab + techniques & no & no code & 81.35 F50 / 71.8 / 54.08 Tab F1 \\
\bottomrule
\end{tabularx}
\end{table}

\subsection{Decisions taken, and the evidence for each}

\begin{table}[h]
\centering
\caption{Design decisions in the transcription pipeline (all November 2026).}
\label{tab:decisions}
\small
\begin{tabularx}{\textwidth}{@{}>{\raggedright}p{4.3cm}X>{\raggedright\arraybackslash}p{1.5cm}@{}}
\toprule
Decision & Evidence & Section \\
\midrule
Modular pipeline, not end-to-end & Wins by 0.17--0.23 Tab F1 on electric guitar and mixes & \ref{sec:bench-tab} \\
BS-RoFormer SW as default separator, Demucs as fallback & +0.07 F50 on synthetic mixes, +0.12 on real songs; cost 34\,s against 10\,s per song & \ref{sec:separation}, \ref{sec:bench-moises} \\
Keep separator + CRNN rather than MuScriptor on the mix & Real songs: MuScriptor 0.42 on the mix against 0.71 for v2 on the stem; 1.3$\times$ real time & \ref{sec:bench-moises} \\
Beat This! final0 over small0 & Accompaniment CMLt 0.80 $\to$ 0.92; bar lines 35 $\to$ 41 of 60 & \ref{sec:rhythm} \\
v2 electric CRNN as an option, not the default & ``Metal High Gain'' 0.43 $\to$ 0.69, but EGDB has no licence & \ref{sec:notes} \\
Learned fingering tagger when installed & String accuracy +0.07 solo, +0.13 EGDB; final-tab 68 $\to$ 78\% & \ref{sec:fingering} \\
No hand-written octave correction for distortion & Made clean and crunch worse, metal no better & \ref{sec:notes} \\
Pan split for several guitars & Lead F1 0.40 $\to$ 0.85 in two-guitar mixes; single guitars unaffected & \ref{sec:separation} \\
Train in TF32, not bf16 & bf16 crashed (NaN) at step 1100 & \ref{sec:notes} \\
Separation-aware training (T1) moved down, fingering (T2) up, after BS-RoFormer & Stems cost only about 0.02 F50 against isolated guitar on synthetic mixes & \ref{sec:next} \\
\bottomrule
\end{tabularx}
\end{table}

\subsection{Failure modes, collected}

\begin{enumerate}
\item \textbf{Distorted guitar in real mixes} is the largest remaining loss:
  0.56 F50 on MoisesDB distorted guitars against 0.83 on acoustic.
\item \textbf{Two guitars in the same place} of the stereo image stay one
  track.
\item \textbf{Fingering on new players} is limited by training-data diversity.
\item \textbf{Techniques}: most bends, vibrato and slides are missed; palm
  mute, harmonics and dead notes are not detected at all.
\item \textbf{Rhythm}: compound meters, 16th triplets and 32nds; solo-guitar
  recordings without accompaniment.
\item \textbf{Same-pitch notes} cannot be represented by MT3-style
  transcribers (a reason not to adopt them as-is).
\item \textbf{No real-song benchmark}: the use case is not yet measured.
\end{enumerate}

\subsection{Open problems in the literature}

\begin{enumerate}
\item \textbf{Mix-level tablature}: no method outputs string and fret from a
  band recording. Adding string and fret tokens to a mix-trained transcriber
  such as MuScriptor would combine mix robustness with timbre-aware fingering,
  and would also fix the same-pitch limitation (a unison on two strings becomes
  two different tokens).
\item \textbf{Separation for transcription}: separators are trained for SDR,
  transcribers on clean stems; training each on the other's output is untested
  for guitar.
\item \textbf{Part separation}: splitting two guitars of the same kind needs
  spatial cues or a separator conditioned on the part.
\item \textbf{Techniques in context}: learned classifiers evaluated on
  amplified and polyphonic material.
\item \textbf{A public real benchmark}: commercial recordings with aligned
  reference tabs.
\item \textbf{Rhythm notation}: quantisation, tuplets and compound meters are
  done by rules and not evaluated anywhere.
\end{enumerate}

\subsection{What the agents are working on next}
\label{sec:next}

The detailed plan, with data and compute estimates, is in
\file{docs/transcription_roadmap.md}; the tasks are named T1--T5 there.

\begin{nextsteps}[: in order]
\begin{enumerate}
\item \textbf{Real-song benchmark} (Phase 0.3) \planned: collect 30--50 song +
  tab pairs, align, hand-check 20--30, record the current pipeline as baseline.
\item \textbf{T1, separation-aware CRNN} \wip: train on what inference sees.
  The uncommitted script \code{scripts/train/pseudo_label_moisesdb.py} builds
  training data from MoisesDB: the input is a separator's guitar stem of the
  real mix, the target the notes a teacher model (v2) transcribed from the
  song's isolated guitars (\term{pseudo-labels}), split by song into train
  (67\%), validation (8\%) and test (25\%). \code{finetune_crnn.py} gains
  \code{--extra-train} and \code{--extra-val} to mix these items in. The target
  is the 0.56 on distorted guitars in real songs.
\item \textbf{T2, fingering}: request DadaGP for pre-training (more players),
  then condition on audio (CRNN features at each onset, as TART does); retrain
  without EGDB if its authors do not allow distribution \planned.
\item \textbf{T3, per-note technique classifier} on GOAT, IDMT and Guitar-TECHS
  with tone augmentation \planned.
\item \textbf{T4, guitar tab transformer on full mixes}: fine-tune MuScriptor
  with string and fret tokens, small first, then large across the DGX Sparks;
  only after T1/T2 baselines exist \planned.
\item \textbf{T5, two guitars in the same spot}: a query-based or fine-tuned
  RoFormer separator for ``rhythm'' and ``lead''; lowest priority \planned.
\end{enumerate}
Smaller items: try Banquet and AlignBeat; palm mute and harmonics; 6/8,
16th triplets and 32nds; batch transcription of a Spotify playlist via YouTube;
a native (ONNX) separator so Android could separate; a tab editor.
\end{nextsteps}

\subsection{Conclusion}

Guitar tablature transcription moved from frame-wise CNNs on acoustic guitar to
modular systems with learned fingering and techniques, and to generative
end-to-end models; in parallel, large transcribers trained on real mixes now
handle distorted guitar but stop at notes. For transcribing a guitar part from a
song, the practical state of the art is a modular pipeline, and that is what the
app implements. A band-split separator recovers most of what the mix costs, a
note model trained on many tones handles amplified guitar on isolated
recordings, and a learned fingering model fixes most of the Viterbi's trouble
with solo lines. What remains is distorted guitar inside real mixes, the
separation of guitars that share a position, techniques beyond the easy
cases, and, before anything else, a benchmark of real songs with reference tabs
so that progress on the actual task can be measured.


\appendix
\part*{Appendices}
\addcontentsline{toc}{part}{Appendices}
% !TeX root = ../main.tex
\section{Glossary}
\label{app:glossary}

\begin{description}[style=unboxed,leftmargin=0cm,itemsep=3pt]
\item[ADRess] Azimuth discrimination and resynthesis: separating sounds by their
  left/right position in a stereo mix (\S\ref{sec:pan-split}).
\item[Autoregressive] A decoder that writes one token at a time, each depending
  on the previous ones; slow on CPU.
\item[Bar-pointer model] A Viterbi/HMM whose state is ``beat $k$ of a bar of $n$
  beats''; used to decode downbeats and meter.
\item[Bleed] Traces of other instruments left in a separated stem.
\item[CMLt] Beat-tracking continuity metric: only runs of consecutive correct
  beats count.
\item[Comp / solo] GuitarSet's two versions of each excerpt: accompaniment
  (chords) and melody.
\item[CQT] Constant-Q transform: a spectrogram with bins at fixed musical
  intervals.
\item[CRNN] Convolutional network followed by a recurrent layer.
\item[DBN] Dynamic Bayesian network; in beat tracking, the decoder that imposes
  tempo and meter.
\item[DI] Direct input: the clean guitar signal before any amplifier.
\item[Downbeat] First beat of a bar.
\item[F50] Note F1 with a 50\ms{} onset tolerance and the right pitch.
\item[Final-tab F1] Our note F1 at 80\ms{} on the finished, quantised tab.
\item[Fine-tuning] Continuing to train an existing model on new data.
\item[GP5] Guitar Pro 5 file format, the app's tab output.
\item[Held out] Not used in training or for choosing settings.
\item[Hexaphonic pickup] A pickup with one output per string; gives string
  labels.
\item[Intermodulation] Sum and difference frequencies created by distortion.
\item[Mask] Per time-frequency-bin share of energy assigned to a source.
\item[mir\_eval] The standard Python library of music-information-retrieval
  metrics.
\item[MSST] Music-Source-Separation-Training, the framework that runs
  BS-RoFormer.
\item[ONNX / tract] A file format for neural networks / the Rust library the app
  runs them with.
\item[Onset, offset] Start and end of a note.
\item[Pan] Position of a sound between the left and right channels.
\item[Posteriorgram] A model's per-frame, per-pitch probabilities.
\item[Pseudo-label] A label produced by a model (the \term{teacher}) instead of
  a human.
\item[Quantisation] Snapping note times onto a grid of beat subdivisions.
\item[SDR] Signal-to-distortion ratio, the usual separation metric, in dB.
\item[Sidecar] The \code{.gvtab.json} file next to a transcribed tab (audio
  hash, engines, bar times).
\item[Stem] One instrument's audio, as produced by a separator or a multitrack.
\item[String accuracy] Share of correct notes placed on the reference string.
\item[Tab F1] Note F1 that also requires the reference string.
\item[Tagger] A model that labels each input item (here: scores the six strings
  per note) instead of generating a sequence.
\item[Tatum] The smallest regular time grid in a piece.
\item[Viterbi] Dynamic-programming search for the cheapest sequence of choices
  (\S\ref{sec:viterbi}).
\item[Zero-shot] Tested on data of a kind never seen in training.
\end{description}

% !TeX root = ../main.tex
\section{Code map and how to reproduce the numbers}
\label{app:code}

\subsection{Source files}

\begin{center}
\small
\begin{tabularx}{\textwidth}{@{}>{\raggedright}p{6.2cm}X@{}}
\toprule
File (under \file{guitar_core/src/} unless noted) & Role and section \\
\midrule
\file{transcribe/mod.rs} & \code{transcribe::run}, stages, progress, sidecar (\S\ref{sec:app}) \\
\file{transcribe/separation.rs} & finds and runs BS-RoFormer or Demucs (\S\ref{sec:separation}) \\
\file{transcribe/stereo.rs} & pan split into guitar parts (\S\ref{sec:pan-split}) \\
\file{transcribe/guitar_crnn.rs} & GAPS / v2 CRNN, windows, post-processing (\S\ref{sec:notes}) \\
\file{transcribe/basic_pitch.rs} & Basic Pitch notes and contour (\S\ref{sec:notes}) \\
\file{transcribe/beat_net.rs}, \file{transcribe/beats.rs} & Beat This!, bar-pointer decoding, DSP fallback (\S\ref{sec:rhythm}) \\
\file{transcribe/arrange.rs} & quantisation, Viterbi fingering, tuning, notation (\S\ref{sec:rhythm}, \S\ref{sec:fingering}) \\
\file{transcribe/fingering_net.rs} & learned fingering tagger (\S\ref{sec:fingering}) \\
\file{transcribe/techniques.rs} & technique rules (\S\ref{sec:techniques}) \\
\file{transcribe/youtube.rs} & yt-dlp search and download (\S\ref{sec:app}) \\
\file{transcribe/eval.rs} & Rust end-to-end benchmarks (\S\ref{sec:benchmark}) \\
\file{shared/tab/export.rs}, \file{shared/tab/repository.rs} & GP5 export, tab repository and sidecars \\
\file{frontend/ui/transcribe.rs} & the Transcribe form \\
\file{guitar_desktop/src/transcribe.rs} & batch CLI \\
\file{scripts/train/} (repository root) & data preparation, fine-tuning, Python benchmarks, tab alignment \\
\file{scripts/train/fingering/} & fingering tagger training \\
\bottomrule
\end{tabularx}
\end{center}

\subsection{Commands}

All Rust benchmarks are ignored tests; run them at workspace level and in
release mode:
\begin{verbatim}
cargo test --release --workspace --exclude guitar_android <name> \
    -- --ignored --nocapture
\end{verbatim}
\begin{center}
\small
\begin{tabularx}{\textwidth}{@{}>{\raggedright}p{3.4cm}X@{}}
\toprule
Test & Inputs (environment variables) \\
\midrule
\code{eval_guitarset} & GuitarSet path; \code{GVA_EVAL_PLAYERS=03,04,05}; \cmd{GVA_EVAL_TONE="Metal High Gain"}; \code{GVA_EVAL_GUITAR_MODEL=electric}; \code{GVA_EVAL_FINGERING=learned} \\
\code{eval_transcription} & Guitar Pro files rendered with a SoundFont \\
\code{eval_fingering} & true notes; \code{GVA_FIT_FINGERING=1} refits the Viterbi weights \\
\code{eval_guitar_pairs} & two GuitarSet guitars per stereo mix; \code{GVA_SEPARATOR}; compare with \code{GVA_SPLIT_GUITARS=0} \\
\code{eval_manifest} & \code{GVA_EVAL_MANIFEST=<tsv>} from \code{prepare_data.py} / \code{make_mixes.py} \\
\code{eval_real} & \code{GVA_EVAL_REAL}: GP8 files with audio, or \code{eval_real.tsv} from \code{align_tabs.py} \\
\bottomrule
\end{tabularx}
\end{center}
\code{GVA_EVAL_CACHE=<dir>} caches each stem's notes so that benchmarks of the
later stages skip the models. The Python stage benchmarks
(\code{bench_notes.py}, \code{bench_separation.py}, \code{bench_beats.py},
\code{bench_tab.py}, \code{bench_moisesdb.py}) and the training steps are
documented with full command lines in \file{scripts/train/README.md}.

\begin{caution}[: rule for anyone changing thresholds]
Run the relevant benchmarks before and after the change, on the same items,
and keep a change only if it wins on the sets it targets without losing on the
others. Never fit anything on GuitarSet players 03--05.
\end{caution}


\bibliographystyle{plainnat}
\bibliography{references}

\end{document}
